% draft3.tex -- Coherence pass of draft2.tex
% Restructured Discussion, compressed f_pi, honest parameter counting
\documentclass[12pt,a4paper]{article}

\usepackage{amsmath,amssymb,amsfonts}
\usepackage{graphicx}
\usepackage{hyperref}
\usepackage{booktabs}
\usepackage{multirow}
\usepackage{geometry}
\geometry{margin=2.5cm}

% Shortcuts
\newcommand{\SU}{\mathrm{SU}}
\newcommand{\UU}{\mathrm{U}}
\newcommand{\SO}{\mathrm{SO}}
\newcommand{\Nf}{N_f}
\newcommand{\Nc}{N_c}
\newcommand{\Tr}{\mathrm{Tr}}
\newcommand{\disc}{\mathrm{disc}}
\newcommand{\adj}{\mathrm{adj}}
\newcommand{\diag}{\mathrm{diag}}
\newcommand{\keV}{\,\mathrm{keV}}
\newcommand{\MeV}{\,\mathrm{MeV}}
\newcommand{\GeV}{\,\mathrm{GeV}}
\newcommand{\Sp}{\mathrm{Sp}}
\newcommand{\TeV}{\,\mathrm{TeV}}
\newcommand{\ie}{i.e.\@}
\newcommand{\eg}{e.g.\@}

\title{Three generations from $\SU(3)$ confinement\\
  and $\UU(3)$ family symmetry}

\author{A.\ Rivero\thanks{Corresponding author. \texttt{rivero@unizar.es}}}
\date{Draft --- \today}

\begin{document}
\maketitle

\begin{abstract}
  We study the flavour sector of the Standard Model using two
  ingredients: the moduli space of $\mathcal{N}=1$ $\SU(3)$ SQCD with
  three flavours, and a $\UU(3)$ family symmetry with canonical
  normalization. The Seiberg confined description gives a meson moduli
  space with a $\mathbb{Z}_2$ involution $M \to \Lambda^4 M^{-1}$ that
  exchanges the $\mathbf{15}$ and $\overline{\mathbf{15}}$ of
  $\SU(5)_{\rm flavour}$. The $\UU(3)$ structure determines the mass
  deformations through a charge ansatz $m_i = k(z_0 + z_i)^2$, with
  $z_0 = 1/\sqrt{6}$ fixed by group theory. The self-dual charge
  spectrum has algebraic discriminant~21, which equals the dynamical
  discriminant of the family Higgs potential only for $N=3$
  generations. Extension to $\Nf=4$ true s-confinement recovers the
  Pati--Salam decomposition from the baryon sector, with the top quark
  as the decoupling fourth flavour. The effective Lagrangian involves a
  single meson field: both electroweak Higgs doublets emerge from~$M$
  and its adjugate $\adj(M)$, which is holomorphic in~$M$ alone,
  constraining the ratio $v_u/v_d$ at the superpotential level
  (subject to K\"ahler corrections).
  Electroweak symmetry breaking necessarily breaks SUSY:
  the $F$-term equations are overconstrained, and the
  SUSY-breaking auxiliary fields \emph{are} the lepton masses,
  $|F_{2j}| = m_j$.
  The family confinement
  scale $\Lambda = f_\pi\sqrt{3}/2$ relates the electroweak VEV to
  the lepton mass hierarchy and the K\"ahler normalization of the
  meson condensate.
  The framework has four continuous inputs---a scale~$k$, an
  angle~$\alpha$ (both fixed by $m_e$ and $m_\mu$),
  $m_d$, and~$m_s$---together with a set of
  structural assumptions (Table~\ref{tab:status}).
  From these, it is consistent with thirteen observables including
  the predicted~$m_\tau$ (0.001\% in~$K$), four quark masses (1--4\%, with $m_b$
  sensitive to the $m_d$ input), the
  Cabibbo angle~(0.01\%), $|V_{cb}| = 1/(5R)$~(1\%),
  $|V_{ub}| = 2\sqrt{2}/(7R^3)$~(0.2\%), and
  $\delta_{CP} = \arctan\sqrt{R}$~($0.2\sigma$), where
  $R = (5+\sqrt{21})/2$.
  The inverse Koide ratio for down quarks,
  $K^{-1}(d,s,b)$, agrees to~$0.25\%$ at its conventional
  evaluation scale ($m_{d,s}$ at $2\GeV$, $m_b$ at~$m_b$) and
  improves to~$0.13\%$ at~$M_Z$.
\end{abstract}


%======================================================================
\section{Introduction}
\label{sec:intro}
%======================================================================

The Standard Model contains 13~free parameters in the flavour
sector---six quark masses, three charged-lepton masses, three CKM
angles, and a CP phase---with no known relations among them.

In 1983, Koide~\cite{Koide:1983qe} constructed a preon model
aiming to relate the Cabibbo angle to the mass ratios between
generations. The model produced, as a byproduct, a mass matrix
ansatz $m_i \propto (a + b_i)^2$ where $a$ is a universal constant
and the $b_i$ are signed quantities summing to zero, with the
constraint that the average of the $b_i^2$ equals~$a^2$.
This predicted the tau mass as $m_\tau = 1776.97\MeV$---confirmed
years later by experiment. The original preon model was subsequently
ruled out by the absence of compositeness signals at high-energy
colliders, but the mass matrix relation survived as an empirical
fact in search of a theoretical home.
The formula is a statement about the mass matrix, not about
square roots: the $b_i$ can be positive or negative, and it is
only their \emph{squares} that must be positive (as befits
physical masses).

Seiberg's exact results for $\mathcal{N}=1$
SQCD~\cite{Seiberg:1994bz,Intriligator:1995au} provide a
controlled setting in which the infrared dynamics of a strongly coupled
gauge theory is described by gauge-invariant composites (mesons and
baryons) subject to an exact quantum constraint. For $\SU(3)$ with
$\Nf=3$ flavours, this constraint defines a moduli space with a
non-trivial involution. In this work we show that, when the mass
deformations of this moduli space are parameterized by the charges of a
$\UU(3)$ family symmetry---where $a = z_0$ is the $\UU(1)$
normalization and the $b_i = z_i$ are the $\SU(3)$ Cartan
charges---the Koide relation becomes an algebraic identity:
\begin{itemize}
  \item The charged-lepton Koide ratio $K \equiv \sum m_i / (\sum
    \sqrt{m_i})^2 = 2/3$ to 0.001\%~\cite{Koide:1983qe},
    evaluated at pole masses (here all $z_0 + z_i > 0$, so
    $\sqrt{m_i} = \sqrt{k}\,(z_0{+}z_i)$ with the positive root).
  \item The down-quark masses satisfy $K(1/m_d, 1/m_s, 1/m_b) = 2/3$
    to 0.25\%~\cite{Rivero:2024epjc}, using
    $\overline{\rm MS}$ masses at their conventional scales
    ($m_{d,s}$ at~$2\GeV$, $m_b$ at~$m_b$).
  \item The Cabibbo angle obeys $|V_{us}| \approx
    \sqrt{m_d/m_s}$~\cite{GattoSartoriTonin:1968,Weinberg:1977hb,Fritzsch:1979zq}.
\end{itemize}

The paper is organized as follows and includes pedagogical
explanations of the SUSY techniques for readers more familiar with
the Standard Model than with superfield methods.
Section~\ref{sec:seiberg}
establishes the Seiberg confined description: the meson moduli space,
its involution, and the off-diagonal meson mass matrix, with
introductory material on holomorphy, anomaly matching, and the
classification of SQCD theories.
Section~\ref{sec:u3} introduces the $\UU(3)$ family charges that
parameterize the mass deformations. Section~\ref{sec:involution}
shows how the involution exchanges the charged-lepton and down-quark
charge assignments, identified with the $\mathbf{15} \leftrightarrow
\overline{\mathbf{15}}$ exchange of $\SU(5)_{\rm flavour}$, and
explains the adjugate mechanism.
Section~\ref{sec:uniqueness} proves a uniqueness theorem for $N=3$
generations via discriminant matching.
Section~\ref{sec:nf4} extends to $\Nf=4$ true s-confinement and the
Pati--Salam structure. Section~\ref{sec:numerics} collects
the numerical results. Section~\ref{sec:discussion} discusses
open problems and structural implications.


%======================================================================
\section{$\SU(3)$ SQCD: Confinement, Symmetry Breaking, and the Moduli
  Space}
\label{sec:seiberg}
%======================================================================

\subsection{Why SUSY exact results are trustworthy}

Before presenting the technical setup, we briefly explain why
supersymmetric gauge theories allow exact non-perturbative
statements---a feature absent in ordinary QCD.

In a supersymmetric theory, the superpotential~$W$ is a holomorphic
function of the chiral superfields and the coupling constants.
Holomorphy---the requirement that $W$ depends only on the fields
$\Phi$ and not on their complex conjugates~$\bar\Phi$---is extremely
constraining. (For readers unfamiliar with the term: a holomorphic
function is a ``complex-analytic'' function, one that can be
expanded in a Taylor series in its complex variables. A function
like $|\Phi|^2 = \Phi\bar\Phi$ is \emph{not} holomorphic because it
depends on~$\bar\Phi$; a function like $\Phi^3 + m\Phi$ \emph{is}
holomorphic.) This constraint forbids most quantum corrections.
Seiberg's key insight~\cite{Seiberg:1994bz} was that holomorphy,
combined with symmetries and known limiting behaviour (weak coupling,
large masses), often determines~$W$ uniquely.
The result is an exact effective superpotential valid at
\emph{all} coupling strengths, including the strongly-coupled
regime where perturbation theory fails.

The second pillar is 't~Hooft anomaly matching.
In quantum field theory, an \emph{anomaly} arises when a symmetry
of the classical theory is violated by quantum effects. The most
familiar example is the chiral anomaly in QED, which is responsible
for the decay $\pi^0 \to \gamma\gamma$. 't~Hooft showed that
anomaly coefficients---numbers computed from one-loop triangle diagrams
of global symmetry currents---are \emph{exact}: they receive no
higher-order corrections and are independent of the energy scale.
This means that any proposed low-energy description of a theory
must reproduce the same anomaly coefficients as the UV quarks and
gluons. This is a powerful consistency check: most candidate
effective theories fail it.
For $\SU(\Nc)$ with $\Nf$ flavours, these conditions---together
with holomorphy---uniquely fix the effective superpotential
in terms of composite fields (mesons and baryons), up to an
overall normalization set by the dynamical scale~$\Lambda$.

The physical picture is analogous to QCD: the strong $\SU(3)_F$
interaction confines the ``preons'' ($Q$,~$\tilde Q$) into
colour-singlet bound states. In QCD, the quark bilinear
$\bar q_i q_j$ forms a $3\times3$ matrix of bound states---the
scalar mesons ($\sigma$,~$\pi$, $K$, $\eta$, $a_0$, \ldots). In our
SUSY theory, the analogous bilinear $Q^i\tilde Q_j$ forms a
$3\times3$ matrix of meson \emph{superfields}~$M^i{}_j$, each
containing a complex scalar and a Weyl fermion (the ``mesino'').
Similarly, the colour-antisymmetric combination of three quarks
(the baryon) has a SUSY counterpart
$B = \epsilon_{abc} Q^{a1} Q^{b2} Q^{c3}$.

The crucial difference from QCD is that supersymmetry and holomorphy
allow \emph{exact} control of the resulting low-energy theory.
In QCD, the chiral Lagrangian for pions is only an
approximation---corrections are incalculable at strong coupling.
In SUSY QCD, the superpotential of the confined theory is exact
to all orders.  This is what makes the present framework
non-trivially predictive: the relations between
meson VEVs and mass deformations (eq.~\ref{eq:vacuum} below)
are not approximations but exact consequences of holomorphy and symmetry.

\paragraph{Why $\Nf = \Nc$ is special: the quantum-modified moduli space.}
Seiberg's analysis~\cite{Seiberg:1994bz,Intriligator:1995au}
classifies SQCD theories by the ratio $\Nf/\Nc$:
\begin{itemize}
  \item $\Nf > 3\Nc$: the theory is not asymptotically free.
  \item $\frac{3}{2}\Nc < \Nf < 3\Nc$: the theory flows to
    an interacting conformal fixed point (``conformal window'').
  \item $\Nc + 1 \leq \Nf \leq \frac{3}{2}\Nc$: the theory confines
    and can be described by dual variables (Seiberg duality).
  \item $\Nf = \Nc + 1$: \emph{s-confinement}. The confined theory
    has a smooth moduli space with a SUSY-preserving origin.
    The effective superpotential is a polynomial in the composites.
  \item $\Nf = \Nc$ (our case): the confined theory has a
    \emph{quantum-modified} moduli space. The classical moduli space
    (where $\det M = 0$ is allowed) is deformed by
    instantons---non-perturbative tunnelling events between
    topologically inequivalent gauge-field configurations---to
    $\det M - B\tilde B = \Lambda^{2\Nc}$, removing the origin.
    This is a genuinely quantum-mechanical effect with no classical
    analogue.
\end{itemize}
The $\Nf = \Nc$ case is the borderline between the s-confining
window and the ADS (Affleck--Dine--Seiberg)
regime~\cite{AffleckDineSeiberg:1984}, where a
dynamical superpotential $W \propto 1/\det M$ is generated and
no stable vacuum exists.  At $\Nf = \Nc$, the ADS superpotential
is replaced by a constraint---a qualitative change that occurs
precisely at the threshold where the baryon becomes an independent
degree of freedom (for $\Nf < \Nc$, baryons cannot be formed).

The reader unfamiliar with SUSY gauge theories may consult
the reviews~\cite{Intriligator:1995au,Seiberg:1994bz} for
pedagogical treatments; we now summarise the essential results.

\subsection{The UV theory and its symmetries}

Consider $\mathcal{N}=1$ $\SU(3)$ SQCD with $\Nf = 3$ flavours of
chiral superfields $Q^i$, $\tilde{Q}_j$
($i,j = 1,2,3$)~\cite{Seiberg:1994bz,Intriligator:1995au}.
(A \emph{chiral superfield} is the basic building block of
$\mathcal{N}=1$ SUSY: it packages a complex scalar~$\phi$, a Weyl
fermion~$\psi$, and an auxiliary field~$F$ into a single object.
The scalar and fermion have the same gauge quantum numbers---they
are ``superpartners.'' The field $Q^i$ is the SUSY analogue of
a quark: it transforms as a fundamental $\mathbf{3}$ of the gauge
group $\SU(3)$, and carries a flavour index $i = 1,2,3$. The
conjugate field $\tilde Q_j$ transforms as $\bar{\mathbf{3}}$.)

The classical global symmetry is
\begin{equation}
  G_{\rm cl} = \SU(3)_L \times \SU(3)_R \times \UU(1)_B
    \times \UU(1)_A \times \UU(1)_R\,.
  \label{eq:Gcl}
\end{equation}
The axial $\UU(1)_A$ is anomalous: although it is a symmetry of the
classical Lagrangian, quantum effects (instantons) break it. This is
exactly analogous to the axial anomaly in QCD, which is responsible
for the large mass of the $\eta'$ meson.\footnote{In QCD, the axial
$\UU(1)_A$ symmetry $q \to e^{i\gamma_5\theta}q$ is broken by the
triangle anomaly to $\mathbb{Z}_{2\Nf}$. The same mechanism operates
here.} The surviving non-anomalous symmetry includes the
$R$-symmetry---a continuous symmetry unique to SUSY theories in which
different components of a superfield carry different charges.
(The $R$-charge is the SUSY analogue of chirality: the scalar
component of a superfield carries $R$-charge~$r$, while the fermion
component carries $R$-charge $r{-}1$. The superpotential must have
total $R$-charge~2, which severely constrains its form.)
With the standard assignment $R_Q = R_{\tilde{Q}} = 0$ and
$R_\Lambda = 2/3$, the $R$-symmetry constrains
the effective superpotential.

\subsection{Confined spectrum}

The low-energy degrees of freedom are gauge-invariant composites:
\begin{itemize}
  \item 9~mesons: $M^i{}_j = Q^i \tilde{Q}_j$
    \quad ($\mathbf{3}_L \otimes \bar{\mathbf{3}}_R$),
  \item a baryon: $B = \epsilon_{abc}\,\epsilon_{ijk}\,Q^{ai}Q^{bj}Q^{ck}$
    \quad ($\mathbf{1}_L$, $B = +1$),
  \item an antibaryon: $\tilde{B} = \epsilon^{abc}\,\epsilon^{ijk}\,
    \tilde{Q}_{ai}\tilde{Q}_{bj}\tilde{Q}_{ck}$
    \quad ($\mathbf{1}_R$, $B = -1$).
\end{itemize}
Each meson superfield $M^i{}_j$ contains a complex scalar and a Weyl
fermion (the ``mesino'') sharing the same quantum numbers. The 9~mesino
fermions transform as $(\mathbf{3}, \bar{\mathbf{3}})$ under
$\SU(3)_L \times \SU(3)_R$: 3~diagonal (neutral under
$\SU(3)_{\rm diag}$) and 6~off-diagonal (charged). This counting---9
Weyl fermions---matches the minimal SM lepton sector
($L_i + e^c_i$, no $\nu_R$)~\cite{ArkaniHamed:1998wc}, a point we return to in the discussion.

These composites satisfy the exact quantum constraint
\begin{equation}
  \det M - B\tilde{B} = \Lambda^6\,,
  \label{eq:constraint}
\end{equation}
where $\Lambda$ is the dynamical scale. Equation~\eqref{eq:constraint}
is enforced by a Lagrange multiplier superfield~$X$ through the
effective superpotential
\begin{equation}
  W_{\rm dyn} = X\bigl(\det M - B\tilde{B} - \Lambda^6\bigr).
  \label{eq:Wdyn}
\end{equation}
This is the quantum-modified moduli space of the $\Nf = \Nc$
theory~\cite{Seiberg:1994bz}. (The $\Nf = \Nc$ theory is \emph{not}
s-confining in the strict sense of~\cite{CsakiSchmaltzSkiba:1997}:
the origin is removed and the moduli space is modified, whereas true
s-confinement at $\Nf = \Nc + 1$ has a smooth origin. See
Section~\ref{sec:nf4}.)

\paragraph{Physical meaning of the quantum constraint.}
The constraint $\det M = \Lambda^6$ on the mesonic branch says
that the product of the three diagonal meson VEVs is \emph{fixed}
by the strong dynamics:
\begin{equation}
  \langle M_1\rangle\,\langle M_2\rangle\,\langle M_3\rangle
  = \Lambda^6\,.
  \label{eq:detfixed}
\end{equation}
This is the SUSY analogue of the QCD chiral condensate
$\langle\bar qq\rangle \sim \Lambda_{\rm QCD}^3$, but elevated
from an approximate order parameter to an exact constraint.
While in QCD the value of the chiral condensate is only known
numerically (from lattice calculations or sum rules), the SUSY
constraint is exact.

The constraint has a dramatic physical consequence: if one meson
VEV is made small (by giving the corresponding preon a large mass),
the others must grow to compensate. This is the origin of
the seesaw-like relation $\langle M_i\rangle \propto 1/m_i$ in
eq.~\eqref{eq:vacuum}---the same structure that produces the
Koide mass hierarchy from a hierarchy of input charges.
Removing the origin ($\det M = 0$ is forbidden) means
that the confined theory has no symmetric vacuum: at least one
meson must condense, and chiral symmetry is always broken.

\subsection{Branches and symmetry breaking}

\paragraph{Moduli space.}
In SUSY theories, families of degenerate vacua parameterised by
scalar VEVs form a \emph{moduli space}.
For $\Nf = \Nc$, it is the constraint surface
$\det M - B\tilde B = \Lambda^6$: a 5-complex-dimensional manifold.

The moduli space has two branches:
\begin{itemize}
  \item \textbf{Mesonic branch} ($B = \tilde{B} = 0$):
    $\det M = \Lambda^6$, so at least one meson VEV is nonzero. The
    diagonal vacuum $\langle M \rangle = \diag(\Lambda^2\omega_1,\,
    \Lambda^2\omega_2,\, \Lambda^2\omega_3)$ with
    $\omega_1\omega_2\omega_3 = 1$ breaks
    $\SU(3)_L \times \SU(3)_R \to \SU(3)_V$ (or further, depending on
    the $\omega_i$).
  \item \textbf{Baryonic branch} ($\det M > \Lambda^6$): $B\tilde{B}
    \neq 0$, breaking $\UU(1)_B$. We do not use this branch, but
    it exists as an integral part of the moduli space.
\end{itemize}

\subsection{F-term equations and the mass-deformed vacuum}
\label{sec:fterms}

We now add explicit masses for the preons. In a SUSY theory, mass
terms live in the superpotential: $W_{\rm tree} = m_i M^i{}_i$
(with $m_i > 0$). Physically, giving a preon a mass~$m_i$ means
giving both the scalar and the fermion inside the superfield~$Q^i$
the same mass (SUSY enforces this).  The mass deformation explicitly
breaks the flavour symmetry: different generations get different
masses, just as in the Standard Model.

The total superpotential is
\begin{equation}
  W = X\bigl(\det M - B\tilde{B} - \Lambda^6\bigr) + m_i M^i{}_i\,.
  \label{eq:Wtot}
\end{equation}
To find the vacuum, we set all $F$-terms to zero (which minimises
the scalar potential; see Appendix~\ref{app:primer}).
The $F$-term equations are:
\begin{align}
  F_X &= \det M - B\tilde{B} - \Lambda^6 = 0\,,
    \label{eq:FX}\\
  F_{M^i{}_j} &= X\,(\text{cof}\,M)^j{}_i + m_i\,\delta^j_i = 0\,,
    \qquad\text{($\text{cof}\,M$ = cofactor matrix $\equiv \adj(M)^T$)}
    \label{eq:FM}\\
  F_B &= -X\,\tilde{B} = 0\,,
    \label{eq:FB}\\
  F_{\tilde{B}} &= -X\,B = 0\,.
    \label{eq:FBt}
\end{align}
From \eqref{eq:FB}--\eqref{eq:FBt}: either $X=0$ (which contradicts
\eqref{eq:FM} since $m_i \neq 0$) or $B = \tilde{B} = 0$. Taking the
latter, we are on the mesonic branch with $\det M = \Lambda^6$.

The off-diagonal components of \eqref{eq:FM} force
$\langle M^i{}_j \rangle = 0$ for $i \neq j$ (since the cofactor
matrix of a diagonal $M$ is diagonal). The diagonal components give
\begin{equation}
  \langle M_i \rangle = \frac{\Lambda^2 P^{1/3}}{m_i}\,,
  \qquad P = m_1 m_2 m_3\,.
  \label{eq:vacuum}
\end{equation}
This is the central equation of the framework. It says:
\emph{heavy preons produce light mesons, and vice versa.}
The meson VEV is \emph{inversely proportional} to the preon mass,
like a seesaw. This inverse relation is not assumed but derived
from the $F$-term equations and the quantum constraint.
It is exact in the SUSY limit and holds non-perturbatively.

\paragraph{Symmetry-breaking pattern.}
The mass deformation breaks $\SU(3)_L \times \SU(3)_R$ explicitly.
In the vacuum \eqref{eq:vacuum}:
\begin{itemize}
  \item If $m_1 = m_2 = m_3$: $\SU(3)_V$ is preserved.
  \item If all $m_i$ are distinct: the global symmetry is broken to
    $\UU(1)^2$.
  \item $\UU(1)_B$ is preserved ($B = \tilde{B} = 0$).
  \item $\UU(1)_R$ is broken by $\langle M_i \rangle \neq 0$.
\end{itemize}
This is a rich scaffold for flavour symmetry breaking, driven entirely
by the mass parameters $m_i$---which in Section~\ref{sec:u3} we fix
using $\UU(3)_F$ charges.

\subsection{The meson mass matrix}

The off-diagonal meson masses are
\begin{equation}
  m_{\rm phys}(M^i{}_j) = \frac{m_i m_j}{P^{1/3}}\,,
  \label{eq:offdiag}
\end{equation}
which are $\Lambda$-independent: they depend only on ratios of the
input masses. The diagonal meson mass matrix, obtained from second
derivatives of $W$ on the constraint surface, is
\begin{equation}
  W_S = P^{1/3}\bigl(\diag(v_i^2) - v \otimes v^T\bigr),
  \qquad v_i = \frac{m_i}{P^{1/3}}\,.
  \label{eq:WS}
\end{equation}
This matrix has the properties:
\begin{enumerate}
  \item $\Tr(W_S) = 0$: the eigenvalues sum to zero.
  \item $(W_S)_{kk} = 0$ for all~$k$: the diagonal entries vanish
    (a consequence of $\partial^2 \det M / \partial M_k^2 = 0$, \ie,
    the multilinearity of the determinant).
  \item Rank~2 (one zero eigenvalue from the constraint).
\end{enumerate}
Property~(2) is the key structural result: the Seiberg dynamics is
\emph{transparent} to the mass parameters. If the input masses $m_i$
satisfy a sum rule, the physical spectrum preserves it at leading order
in the off-diagonal perturbation.

\paragraph{Eigenvalue structure.}
The eigenvalues of $W_S$ at the physical lepton mass scale
sum to zero: $\lambda_1 + \lambda_2 +
\lambda_3 = 0$. This immediately rules out identifying the Seiberg
eigenvalues directly with lepton masses (which are all positive).
Numerically, $\lambda \approx \{+3229, -3229, -0.4\}\MeV$ for
physical lepton mass inputs. The Seiberg dynamics preserves the charge
structure but does not generate the mass spectrum---the mechanism lives
in the family sector.

\paragraph{Summary of the confined framework.}
To summarise the physical picture before proceeding to the charge
ansatz: three ``preons'' ($Q_i$ and $\tilde Q_j$) bind under
a family $\SU(3)_F$ gauge force to form 9~meson superfields
$M^i{}_j$ and a baryon/antibaryon pair. The mesons contain the
Higgs doublets and the lepton fields. The vacuum is determined by
an exact quantum constraint (the product of meson VEVs is fixed)
plus mass deformations $m_i$ that break the flavour symmetry.
Heavy preons produce light mesons: $\langle M_i\rangle \propto
1/m_i$. The off-diagonal meson mass matrix is ``transparent''---it
does not spoil the charge structure. The remaining question is:
\emph{what determines the mass deformations $m_i$?}


%======================================================================
\section{Family Charges from $\UU(3)$}
\label{sec:u3}
%======================================================================

The mass parameters $m_i$ in the tree-level superpotential
$W_{\rm tree} = m_i M^i{}_i$ are free inputs. We now specify them
using a $\UU(3)_F$ family symmetry.

Let the three generations transform as the fundamental
$\mathbf{3}$ of $\UU(3)_F = \SU(3)_F \times \UU(1)$, with all
generators in canonical normalization $\Tr(T_A T_B) =
\frac{1}{2}\delta_{AB}$.

The $\SU(3)$ Cartan charges of the fundamental components satisfy
\begin{equation}
  \sum_{i=1}^3 z_i = 0\,,
  \qquad
  \sum_{i=1}^3 z_i^2 = \frac{1}{2}\,,
  \label{eq:su3constraints}
\end{equation}
and the $\UU(1)$ charge is fixed:
\begin{equation}
  z_0 = \frac{1}{\sqrt{6}}\,.
  \label{eq:z0}
\end{equation}

\paragraph{Physical interpretation of $m_i$.}
The parameter $m_i$ is the tree-level superpotential mass deformation
for the preon bilinear $Q^i\tilde Q_i$:
$W_{\rm tree} = \sum_i m_i\,Q^i\tilde Q_i$.
After confinement, $m_i$ appears simultaneously as: (i)~the mass
deformation that lifts the moduli space; (ii)~the effective Yukawa
coupling that determines fermion hierarchies (Section~\ref{sec:involution});
and (iii)~the determinant factor in the meson VEV via
$\langle M_i\rangle = \Lambda^2 P^{1/3}/m_i$.
As a superpotential coupling, $m_i$ has mass dimension one and is
protected by holomorphy against perturbative corrections
in the $\mathcal{N}=1$ theory.

\paragraph{Charge ansatz.}
We \emph{postulate} a quadratic relation between the mass deformation
and the family quantum numbers:
\begin{equation}
  m_i = k\,(z_0 + z_i)^2\,.
  \label{eq:massansatz}
\end{equation}
This is the norm-squared on the $\UU(3)$ charge lattice---the
Killing-form inner product restricted to the Cartan subalgebra.
The quadratic form is distinguished on three grounds:
(i)~it is the unique $W$-invariant bilinear on the
weight lattice (the Weyl group permutes the $z_i$, leaving
$\sum(z_0 + z_i)^2$ invariant);
(ii)~in heterotic string compactifications on a Narain lattice,
the BPS mass-squared is
$M^2 = |p_L|^2/\alpha'$~\cite{Narain:1986}, where $p_L$ is the
left-moving momentum on the charge lattice; a Wilson line provides
the universal shift $z_0$, the $z_i$ are gauge Cartan charges;
(iii)~in the $\mathcal{N}=2$ attractor mechanism, the BPS mass
at the attractor point is $M^2 = |Z|^2$ with $Z$ linear in
charges~\cite{Ferrara:1995ih}; on the $\UU(3)$ lattice this
gives \eqref{eq:massansatz}.
Note that the standard BPS formula $M^2 \propto Q^2$ yields a
mass-\emph{squared} quadratic in charges, whereas the ansatz
\eqref{eq:massansatz} has a mass (dimension~1) quadratic in charges.
The meson $M^i{}_j = Q^i\tilde{Q}_j$ resolves this: as a
quark--antiquark bilinear it has mass dimension~2, so its
VEV $\langle M_i\rangle$ plays the role of $M^2$, not~$M$.
Physically, the meson is a flavour ``capacitor''---a bound pair
of charges $Q^i$ and $\tilde{Q}_i$ separated at the confinement
scale~$\Lambda_F$.  The quadratic form has a gauge-theoretic origin:
by Gauss's law the chromoelectric field of a static preon with
family charge~$q_i$ is $|\mathbf{E}_i| \propto q_i$; since
the gauge Lagrangian is $\mathcal{L} \supset -(1/4g^2)F_{\mu\nu}^2$,
the field energy is $U_i = (2g^2)^{-1}\!\int\!d^3x\,\mathbf{E}_i^2
\propto q_i^2\,\Lambda_F/g_F^2$,
i.e.\ quadratic in the charge, just as for a classical capacitor
$E = Q^2/(2C)$.
The same result follows from the $\mathcal{N}=1$ effective
Lagrangian for the residual $\UU(1)^2_F$ after confinement.
On the mesonic vacuum with $B = \tilde{B} = 0$, the
static energy functional
$\mathcal{E} = (2g^2)^{-1}(\mathbf{E}^2{+}\mathbf{B}^2)
+ \tfrac{g^2}{2}J_F^2 + \cdots$ contains the $D$-eliminated
family current $J_F = \sum_i q_i\,|\phi_i|^2 + \xi$
(where $\phi_i$ are the scalar meson fields).
For composite states whose spatial profile is generation-independent
(the confinement scale~$\Lambda_F$ sets the tube width for all~$i$)
but whose field strength scales as $\mathbf{E}_i = q_i\,
\mathbf{E}_0$, the integrated energy is
$\Delta M_i = q_i^2\,\mu$ with
$\mu = \!\int\!d^3x\bigl[(2g^2)^{-1}\mathbf{E}_0^2
+ (g^2/2)\,J_0^2\bigr]$.
The overall scale~$\mu = k$ is non-perturbative and depends on
the K\"ahler metric; the $q_i^2$~pattern is fixed by gauge theory.
In the M-theory lift (Appendix~\ref{app:Mtheory}), this
bound pair maps to an M2-brane whose BPS mass is
$T_2 \times \text{Area}$; the quadratic form would follow if
both worldvolume directions scale linearly with
$(z_0{+}z_i)$---a geometric restatement
of the same field-energy picture.

\paragraph{What is excluded---and why the square is unique.}
A linear ansatz $m_i = k'(z_0 + z_i)$ is excluded by
positivity (some $m_i < 0$ for $|z_i| > z_0$). A quartic ansatz
$m_i = k''(z_0+z_i)^4$ does not produce $K = 2/3$.
More generally, consider the monomial family
$m_i = c\,(z_0 + z_i)^n$. The signed Koide ratio
\begin{equation}
  K_n \;=\; \frac{\sum_i (z_0{+}z_i)^n}
    {\bigl(\sum_i (z_0{+}z_i)^{n/2}\bigr)^2}
  \label{eq:Kn}
\end{equation}
is \emph{direction-independent} in the Cartan subalgebra---i.e.\
independent of the sector angle~$\alpha$---if and only if~$n = 2$.
For $n=2$, the numerator is $S_2 = 3z_0^2 + \sum z_i^2 = 1$
and the denominator is $S_1^2 = (3z_0)^2 = 3/2$, both
independent of~$\alpha$; hence $K_2 = 2/3$ identically.
For $n \geq 3$, $S_n$ depends on the cubic and higher Casimirs
$\sum z_i^p$ ($p\geq 3$), which vary with direction.
The cubic Casimir $d_{abc}\,T^a T^b T^c \propto \mathbf{1}$ for
the fundamental, so it shifts~$z_0$ but not the
generation-dependent~$z_i$: the leading correction that could
spoil $K = 2/3$ is the direction-dependent power sum
$S_3 \supset \cos(3\alpha)/(4\sqrt{3})$, suppressed by an additional
insertion of the family-breaking background field.
Mixed invariants of the form $m_i m_j$ for $i\neq j$ are
off-diagonal mass terms forbidden by the diagonal $\UU(1)^3$
unbroken in the Cartan subalgebra.

The uniqueness of $n=2$ has a physical interpretation:
the family charge~$q_i$ is an \emph{amplitude}; the
mass is its \emph{norm-square}.  Every second-order mechanism
that produces mass from a linear precursor---seesaw mixing
$\mu_i \propto q_i$ giving $m_i \propto \mu_i^2/M$,
bilinear composites $M^i{}_i = Q^i\tilde{Q}_i$ giving a VEV
$\propto q_i^2$, capacitor energy $E = Q^2/(2C)$,
or M2-brane area $\propto (\text{distance})^2$---yields $n=2$
and therefore direction-independent $K = 2/3$.

\paragraph{Radiative stability.}
Since $m_i$ is a superpotential coupling, it receives no
perturbative corrections from gauge loops (the non-renormalization
theorem of~\cite{Seiberg:1993vc}). The leading violation arises
from $\UU(3)_F$-breaking soft terms at the scale $m_{\rm soft}$;
writing $\delta m_i \sim m_{\rm soft}\,c_i$, the correction to $K$
is $\delta K \sim (m_{\rm soft}/k)\,(\delta c/c) \sim 10^{-2}$
for $m_{\rm soft} \sim 1\TeV$, $k \sim 80\GeV$.
Thus the $K=2/3$ relation is radiatively stable to per-cent accuracy
whenever the family symmetry breaking scale exceeds the soft mass scale.
We emphasise that the ansatz~\eqref{eq:massansatz} is
\emph{uniquely selected}: it is the only monomial mass formula
for which $K=2/3$ is a group-theoretic identity rather
than a direction-dependent numerical coincidence.
The field-energy argument above provides a gauge-theoretic
derivation of the pattern ($q_i^2$); the overall scale
$k = \mu$ remains a non-perturbative quantity tied to
the K\"ahler metric (Section~\ref{sec:mild}).

It follows from
\eqref{eq:su3constraints}--\eqref{eq:z0} that the Koide ratio
\begin{equation}
  K \;\equiv\; \frac{\sum m_i}{\bigl(\sum\sqrt{m_i}\bigr)^2}
  = \frac{\sum (z_0{+}z_i)^2}{\bigl(\sum |z_0{+}z_i|\bigr)^2}
  = \frac{2}{3}
  \label{eq:K23}
\end{equation}
identically for \emph{any} direction in the Cartan subalgebra.
Note that $\sqrt{m_i}$ in the standard formula means
$|z_0 + z_i|$---the absolute value of a signed charge. When all
charges are positive ($z_0 + z_i > 0$ for each~$i$), the
absolute value is redundant.
For sectors where one charge crosses zero (as in the $(b,c,s)$
bridge tuple of Section~\ref{sec:numerics}), the sign matters:
the Koide ratio is computed with signed charges,
$K = \sum(z_0{+}z_i)^2/(\sum(z_0{+}z_i))^2$, not with positive
square roots.

The proof uses only $\sum z_i = 0$ (tracelessness), $\sum z_i^2 = 1/2$
(Casimir), and $z_0 = 1/\sqrt{6}$ (canonical $\UU(1)$ normalization).
In $\SU(3) \times \UU(1)$, the value of $z_0$ is free; in $\UU(3)$, it
is fixed by the requirement $\Tr(T_0^2) = 1/2$. String constructions
generate $\UU(N)$, not $\SU(N)$, from $N$~D-branes, with the
$\UU(1)$ normalization protected by
holomorphy~\cite{Polchinski:1998rr}.

\paragraph{Angle parameterization.}
The Cartan charges are
\begin{equation}
  z_i = \frac{1}{\sqrt{3}}\,\cos\!\bigl(\alpha + \tfrac{2\pi (i-1)}{3}\bigr),
  \qquad i = 1,2,3\,.
  \label{eq:zi}
\end{equation}
Each mass sector is parameterized by one angle $\alpha$ and one
scale~$k$: two parameters for three masses.

\paragraph{Transparency.}
Combined with property~(2) of $W_S$ (eq.~\ref{eq:WS}), the ratio
$K=2/3$ is preserved to first order through the Seiberg dynamics.
To see this, write the full meson mass matrix as
$\mathcal{M} = \diag(m_i) + W_S$, where $m_i = k(z_0{+}z_i)^2$
are the UV masses and $W_S = P^{1/3}(\diag(v_i^2) - v\otimes v^T)$
is the Seiberg perturbation. Since $(W_S)_{kk} = 0$ (diagonal entries
vanish by det multilinearity), the diagonal of $\mathcal{M}$ is
$m_k$ itself. By standard first-order perturbation theory,
the eigenvalues are $\lambda_k = m_k + O(|W_S^{\rm off}|^2/\Delta m)$,
and $K(\lambda_k) = K(m_k) + O(v^4/m^4) = 2/3 + O(v^4/m^4)$.
Numerically, $|v^2_{\rm max}/m_{\rm min}| \sim 10^{-3}$,
so the correction is negligible. The confining dynamics
neither generates nor destroys the charge structure; it is
\emph{transparent} to the Koide relation.\footnote{This
argument assumes the off-diagonal entries of $W_S$ are perturbative
compared to the mass splittings. For SUSY-broken meson masses the
relevant splittings are the lepton masses, and $W_S^{ij} = -v_iv_j$
with $v_i = m_i/P^{1/3}$; the largest entry is
$|W_S^{23}| \approx 3.2\GeV$ while $m_\tau - m_\mu \approx 1.7\GeV$,
so the perturbative expansion is only marginally valid for the
$(\mu,\tau)$ sector. A non-perturbative argument using the
exact eigenvalue equation $\det(\mathcal{M} - \lambda) = 0$ confirms
the result~\cite{Rivero:2024epjc}.}


%======================================================================
\section{The Involution and the
  $\mathbf{15} \leftrightarrow \overline{\mathbf{15}}$ Exchange}
\label{sec:involution}
%======================================================================

The constraint hypersurface $\det M = \Lambda^6$ has a remarkable
symmetry: it is invariant under the exchange
\begin{equation}
  M_i \;\to\; \frac{\Lambda^4}{M_i}\,,
  \qquad\text{mapping}\qquad
  m_i \;\to\; \tilde{m}_i = \frac{P^{2/3}}{m_i}\,.
  \label{eq:involution}
\end{equation}
This is a $\mathbb{Z}_2$ \emph{involution}---a map that, applied
twice, returns to the starting point: $(M_i \to \Lambda^4/M_i
\to M_i)$. Physically, it exchanges large and small meson VEVs
while preserving the product constraint. Each point on the moduli
space has a ``mirror'' point related to it by this involution;
the self-dual locus ($M_i = \Lambda^4/M_i$, \ie,
$M_i = \Lambda^2$) is the fixed-point set.

Since $K$ is scale-invariant, $K(\tilde{m}) = K(1/m)$: the involution
exchanges the charge assignment of the standard sector (charged leptons,
$K(m)=2/3$) with that of the inverse sector (down quarks,
$K(1/m)=2/3$). This is the key to the lepton--quark connection:
the \emph{same} UV charge structure produces both Koide relations,
related by the moduli-space involution.

\subsection{The adjugate mechanism}

The dynamical origin of this exchange is the \emph{adjugate} of the
meson matrix---a construction that may be unfamiliar to readers
outside linear algebra, but plays a central role in the framework.

For any $n\times n$ matrix~$A$, the adjugate $\adj(A)$ is the
transpose of the cofactor matrix. For a $3\times3$ diagonal
matrix $M = \diag(M_1, M_2, M_3)$, it takes the simple form
\begin{equation}
  \adj(M)_i = \prod_{j \neq i} M_j\,,
  \qquad
  M \cdot \adj(M) = \det(M)\,\mathbf{1}\,.
  \label{eq:adjugate}
\end{equation}
Explicitly: $\adj(M)_1 = M_2 M_3$, $\adj(M)_2 = M_1 M_3$,
$\adj(M)_3 = M_1 M_2$. The key point is that $\adj(M)$ is a
\emph{polynomial} (quadratic) in the meson entries---it involves
only multiplication, no division. This makes it holomorphic: it
is a legitimate superpotential coupling. Yet from the identity
$M \cdot \adj(M) = \det(M)\,\mathbf{1}$, we see that on the
constraint surface $\det M = \Lambda^6$, the adjugate
\emph{behaves} like the matrix inverse:
$\adj(M) = \Lambda^6\,M^{-1}$.

This is the mathematical magic of the framework: the adjugate
provides a holomorphic realisation of the matrix inverse---an
operation that is \emph{not} holomorphic in general ($1/M$ involves
$M$ in the denominator). The quantum constraint ``trades'' the
non-holomorphic inverse for a holomorphic polynomial.

The $F$-term equations of the s-confined superpotential
$W = (\det M - B\tilde B)/\Lambda^3 + m_i M^i{}_i$ yield, for
diagonal $M$:
\begin{equation}
  \frac{\partial W}{\partial M^i{}_i}
  = \frac{\adj(M)_i}{\Lambda^3} + m_i = 0
  \qquad\Longrightarrow\qquad
  \adj(M)_i = -\Lambda^3\,m_i\,.
  \label{eq:Fterm}
\end{equation}
The $B$, $\tilde B$ $F$-terms force $B = \tilde B = 0$ in the massive
vacuum. Thus \emph{the adjugate of the meson matrix is proportional to
the UV mass matrix}, while $M_i \propto 1/m_i$. This is exact in the
supersymmetric limit.

Standard Model fermion masses arise from effective Yukawa couplings
to the confined composites. If the left chiral symmetry is partially
gauged as $\SU(2)_L \times \UU(1)_Y \subset \SU(3)_L$, with the
branching rule $\mathbf{3}_L \to \mathbf{2}_{-1/2} \oplus
\mathbf{1}_{+1}$, then the meson doublet $M^{1,2}{}_i$ has the
quantum numbers of the down-type Higgs: $(1,2)_{-1/2}$. The adjugate
$\adj(M)$ transforms as $(\bar{\mathbf{3}}_L, \mathbf{3}_R)$; since
$\bar{\mathbf{3}}_L \to \mathbf{2}_{+1/2} \oplus
\mathbf{1}_{-1}$, its doublet component has $(1,2)_{+1/2}$---the
up-type Higgs. Both electroweak Higgs doublets thus emerge from a
single meson matrix~$M$, with $\adj(M)$ a holomorphic polynomial in~$M$
(not an independent chiral superfield). At the superpotential level,
this constrains $\tan\beta$: $\langle\adj(M)_i\rangle/\langle M_i\rangle
= \Lambda^6/\langle M_i\rangle^2$ is fixed by the quantum constraint.

In the Standard Model, fermion masses arise from Yukawa couplings
$y\,\bar\psi\,H\,\psi$ between fermions and the Higgs.
In the present framework, the Higgs is the meson doublet $M^{1,2}$,
and there are two natural holomorphic Yukawa operators at the
confinement scale---one using $M$ directly, the other using its
adjugate:
\begin{align}
  W_{\rm down} &= y_d\,q_L\,d^c\,\frac{M^{1,2}}{\Lambda}
  + y_\ell\,L\,e^c\,M^{1,2}\,,
  &\quad [H_d\text{-type},\; Y{=}-\tfrac{1}{2}]\,,
  \label{eq:W1} \\
  W_{\rm up} &= y_u\,q_L\,u^c\,\frac{\adj(M)^{1,2}}{\Lambda^3}\,,
  &\quad [H_u\text{-type},\; Y{=}+\tfrac{1}{2}]\,.
  \label{eq:W2}
\end{align}
The $W_{\rm down}$ operator gives masses $\propto \langle M_i
\rangle \propto 1/m_i$ (``standard'' masses, large for light
preons); the $W_{\rm up}$ operator gives masses $\propto
\langle\adj(M)_i\rangle \propto m_i$ (``inverse'' masses, large
for heavy preons). This is why the lepton and down-quark sectors
obey \emph{different} Koide relations: they couple to different
operators of the same meson field.
All operators conserve hypercharge.
The mass parameter $m_i$ in $W_{\rm tree} = m_i M^i{}_i$ serves three
simultaneous roles: (i)~the UV mass deformation of $\SU(3)_F$ SQCD,
(ii)~the effective Yukawa coupling between preons $Q_i$ and $\tilde Q_i$,
and (iii)~the determinant of the meson VEV via
$\langle M_i\rangle = \Lambda^2 P^{1/3}/m_i$. There are no independent
Yukawa parameters: the hierarchy of SM masses \emph{is} the hierarchy of
family charges $z_i$.

Species coupling through $M$ acquire physical
masses $m_{\rm phys} \propto M_i \propto 1/m_i$; species coupling
through $\adj(M)$ acquire $m_{\rm phys} \propto \adj(M)_i
\propto m_i$.

Charged leptons, being the mesino components of $M^{1,2}{}_i$ itself,
have $m_\ell \propto 1/m$ and therefore $K(m_\ell) = K(1/m) = 2/3$
(standard Koide). Down quarks, as elementary fields, couple to $M$
through partial compositeness~\cite{Kaplan:1991dc} with
generation-dependent mixing $\lambda_i \propto m_i^{\gamma_*}$.
If $\gamma_* \geq 1$, the net mass scales as
$m_d \propto m_i^{2\gamma_*-1}$, and the physical Koide parameter
depends on the exponent. In particular, $\gamma_* = 1$ gives
$m_d \propto m_i$, producing inverse Koide:
\begin{equation}
  K(m_\ell) = K(1/m_d) = K(1/m) = \frac{2}{3}\,,
  \label{eq:koide_duality}
\end{equation}
unifying the two empirical relations as a single UV condition.

After soft SUSY breaking, the meson VEVs shift from the
$F$-flat locus. The zero-diagonal property $(W_S)_{kk} = 0$ of the
off-diagonal meson mass matrix (Section~\ref{sec:seiberg}) ensures
that the Koide condition is preserved to first order in
$m_{\rm soft}/\Lambda$: the meson sector is \emph{transparent} to the
family-sector structure.

\subsection{SU(5) assignment and arithmetic}

In the $\SU(5)$ flavour decomposition of~\cite{Rivero:2024epjc}, the
meson $M^{ij} \in \mathrm{Sym}^2(5) = \mathbf{15}$ and the adjugate
$\adj(M)_{ij} \in \mathrm{Sym}^2(\bar 5) =
\overline{\mathbf{15}}$. The involution \eqref{eq:involution} maps
$\mathbf{15} \leftrightarrow \overline{\mathbf{15}}$, with the
adjoint $\mathbf{24}$ fixed. Standard and inverse charge assignments
are related by charge conjugation on the scalar composites.
This generalises the Georgi--Jarlskog~\cite{GeorgiJarlskog:1979}
mass relations between quarks and leptons from a GUT boundary condition
to a moduli-space involution.

The involution has an arithmetic interpretation: $R = (5+\sqrt{21})/2$
is the fundamental unit of the ring of integers
$\mathbb{Z}[(1+\sqrt{21})/2]$ of the real quadratic field
$\mathbb{Q}(\sqrt{21})$, with $\mathrm{Norm}(R) = R\cdot R' = 1$.
The map $R \to 1/R = R' = (5-\sqrt{21})/2$ is the non-trivial element
of $\mathrm{Gal}(\mathbb{Q}(\sqrt{21})/\mathbb{Q})$: the moduli-space
involution \emph{is} the Galois conjugation $\sqrt{21}\to -\sqrt{21}$.


%======================================================================
\section{Self-Duality, the Discriminant 21, and Uniqueness of $N=3$}
\label{sec:uniqueness}
%======================================================================

This section addresses the question: \emph{why three generations?}
The Standard Model works for any number of generations
(anomaly cancellation requires each generation to be complete, but
does not fix their number). We show that the charge structure of
the preceding sections uniquely selects $N = 3$.

\subsection{The self-dual charge spectrum}

A natural question to ask about the involution of
Section~\ref{sec:involution} is: what charge spectrum is
\emph{self-dual}---invariant under $m_i \to P^{2/3}/m_i$?
On the self-dual locus, the charges form a geometric progression
$Z_1 : Z_2 : Z_3 = 1 : R : R^2$---the same structure that arises
from periodic solutions of Schwinger--Dyson gap equations, where
Blumhofer and Hutter~\cite{BlumhoferHutter:1997} found geometric
mass ratios $\sim e^\pi$ between generations without free parameters.
Imposing both $K(m)=2/3$
and $K(1/m)=2/3$ is algebraically equivalent to requiring
\begin{equation}
  3(1 + R^2 + R^4) = 2(1 + R + R^2)^2\,,
  \label{eq:koide_sd}
\end{equation}
which, using $1{+}R{+}R^2 = 6R$ if $R^2 = 5R{-}1$, reduces to
$72R^2 = 72R^2$ identically. The self-dual Koide identity
\eqref{eq:koide_sd} is therefore equivalent to the minimal polynomial
\begin{equation}
  t^2 - 5t + 1 = 0\,,
  \qquad
  R = \frac{5+\sqrt{21}}{2}\,,
  \qquad
  \disc = 21\,.
  \label{eq:selfdual}
\end{equation}
The self-dual angle follows from the geometric-progression
condition.  Writing $Z_i = (1 + \sqrt{2}\,c_i)/\sqrt{6}$ with
$c_i = \cos(\alpha + 2\pi(i{-}1)/3)$, the condition $Z_2^2 = Z_1 Z_3$
expands to
\begin{equation}
  3\sqrt{2}\,c_2 + 2c_2^2 - 2\,c_1 c_3 = 0\,.
  \label{eq:geom_prog}
\end{equation}
The combination $2c_2^2 - 2c_1c_3$ is evaluated by
product-to-sum identities (using $c_1 c_3 =
[-\tfrac{1}{2}+\cos(2\alpha{+}\tfrac{2\pi}{3})]/2$ and
$c_2^2 = [1+\cos(2\alpha{+}\tfrac{2\pi}{3})]/2$):
it equals $\frac{3}{2}$ independently of~$\alpha$.
Hence~\eqref{eq:geom_prog} reduces to
$c_2 = -1/(2\sqrt{2})$.
Since $c_2 = \cos(\alpha - 2\pi/3) = \sin(\alpha - \pi/6)$,
we obtain $\sin(\alpha - \pi/6) = -1/(2\sqrt{2})$, whence
$\cos(\alpha - \pi/6) = \sqrt{7/8}$ and
$\tan(\pi/6 - \alpha) = 1/\sqrt{7}$.  Therefore
\begin{equation}
  \alpha_{\rm sd} = \frac{\pi}{6} - \arctan\frac{1}{\sqrt{7}}\,,
  \qquad 7 = \frac{D}{N} = \frac{21}{3}\,.
  \label{eq:alpha_sd}
\end{equation}
The structural form
$\alpha_{\rm sd} = \pi/(2N) - \arctan\sqrt{N/D}$ links the self-dual
angle to the same trio $(\pi, N, D{=}21)$ that controls the uniqueness
theorem and the CP~phase.

\subsection{General-$N$ uniqueness}

The construction of the preceding sections generalises to $N$
generations with a $\UU(N)_F$ family symmetry. Two numbers arise
naturally:
\begin{itemize}
  \item The \textbf{algebraic discriminant} $D_{\rm alg}$: the
    discriminant of the self-dual quadratic $t^2 - (N{+}2)t + 1 = 0$,
    which is $D_{\rm alg} = (N{+}2)^2 - 4$. This characterises the
    number field $\mathbb{Q}(\sqrt{D_{\rm alg}})$ in which the
    self-dual charge ratios live.
  \item The \textbf{dynamical discriminant} $D_{\rm dyn}$: from the
    family Higgs potential at the
    enhanced symmetry point $\lambda_0 = \lambda_1 = \lambda_{\rm eff}$,
    the coupling ratio $\mu_0^2/\mu_a^2 = (N{+}2)/(2N{+}1)$ defines
    $D_{\rm dyn} = N(2N{+}1)$. This counts the degrees of freedom in
    the family sector.
\end{itemize}
For the charge structure and the dynamics to be compatible, these
two numbers must coincide: $D_{\rm alg} = D_{\rm dyn}$.
This condition reduces to $N^2 - 3N = 0$
(since $(N{+}2)^2 - 4 = N^2 + 4N = N(2N{+}1)$ iff $N^2 = 3N$):
\begin{equation}
  \boxed{N = 0 \;\text{ or }\; N = 3\,.}
  \label{eq:N3}
\end{equation}

\begin{table}[ht]
\centering
\begin{tabular}{ccccc}
\toprule
$N$ & $D_{\rm alg}$ & $D_{\rm dyn}$ & $\dim\adj\,\SO(2N{+}1)$ & Match? \\
\midrule
1 & 5  & 3  & 3  & No \\
2 & 12 & 10 & 10 & No \\
\textbf{3} & \textbf{21} & \textbf{21} & \textbf{21} & \textbf{Yes} \\
4 & 32 & 36 & 36 & No \\
5 & 45 & 55 & 55 & No \\
\bottomrule
\end{tabular}
\caption{Discriminant matching uniquely selects $N=3$, where
$D = 21 = \dim\adj\,\SO(7) = \dim\Sp(6)$. The Lie algebras $B_N$ and $C_N$
have the same dimension $N(2N{+}1)$; the symplectic form
$\dim\Sp(2N)$ counts degrees of freedom in the electric--magnetic
charge lattice of $N$~families.}
\label{tab:N3}
\end{table}

A complementary constraint comes from the sector angle.
The requirement $\alpha < \alpha_{\rm crit} = \pi/(4N)$
(necessary for positive fermion masses) restricts the allowed
range; for the fitted lepton angle $\alpha_{\rm lep} =
0.22222(1)\;\text{rad}$, only $N \geq 3$ admits a viable Koide
spectrum---a criterion independent of discriminant matching.
A third criterion:
$N^2 + 4 = N^2 + N + 1$ only at $N = 3$.
The left-hand side appears in the self-dual discriminant
($D_{\rm alg} = N(N{+}4)$); the right-hand side is
the dimension of the projective plane $\mathbb{P}^2(\mathbb{F}_N)$.
The coincidence is equivalent to $D_{\rm alg} = D_{\rm dyn}$.
A fourth, M-theoretic criterion (Appendix~\ref{app:Mtheory}):
$\binom{3N}{N} = (N{+}1)\cdot N(2N{+}1)$ holds uniquely at $N = 3$,
giving $84 = 4 \times 21$.  The factors are the number of
$C_3$~polarisations in 11d supergravity ($\binom{9}{3} = 84$),
the $\Nf = N_c{+}1 = 4$ Pati--Salam flavours, and
$\dim\Sp(6) = 21$---the charge-lattice discriminant.


%======================================================================
\section{$\Nf = 4$, Pati--Salam Structure, and the Hybrid Scenario}
\label{sec:nf4}
%======================================================================

As noted in Section~\ref{sec:seiberg}, the $\Nf = \Nc$ theory has a
quantum-modified moduli space---the origin is removed, and the
constraint $\det M = \Lambda^6$ is enforced by a Lagrange multiplier.
\emph{True} s-confinement~\cite{Intriligator:1995au,CsakiSchmaltzSkiba:1997}---where
the confined theory has a smooth moduli space with a
SUSY-preserving origin, and the effective superpotential is a
polynomial in the composites without any Lagrange
multiplier---requires $\Nf = \Nc + 1$.
For $\SU(3)$, this gives $\Nf = 4$.
The s-confining superpotential is uniquely determined by
holomorphy, symmetries, and the requirement that it reproduces
the correct anomaly matching coefficients (all $11/11$ independent
't~Hooft conditions).

\paragraph{Global symmetry and gauged subgroup.}
The UV theory has global symmetry
$\SU(4)_L \times \SU(4)_R \times \UU(1)_B \times \UU(1)_R$.
The confined spectrum consists of:
\begin{itemize}
  \item 16~mesons $M^i{}_j = Q^i\tilde{Q}_j$, $i,j = 1,\dots,4$,
    transforming as $(\mathbf{4},\bar{\mathbf{4}})$;
  \item 4~baryons $B_i = \epsilon^{abc}Q_a^j Q_b^k Q_c^l
    \epsilon_{ijkl}$, transforming as $(\bar{\mathbf{4}},\mathbf{1})$;
  \item 4~antibaryons $\tilde{B}^j = \epsilon_{abc}
    \tilde{Q}^a_k\tilde{Q}^b_l\tilde{Q}^c_m\epsilon^{jklm}$,
    transforming as $(\mathbf{1},\mathbf{4})$.
\end{itemize}
The confined superpotential is
$W = (M^i{}_j B_i\tilde{B}^j - \det M)/\Lambda^5$,
with all 't~Hooft anomalies matching between UV and IR
($11/11$ independent conditions, vs.\ $7/11$ for $\Nf = 3$).

The Pati--Salam model~\cite{PatiSalam:1974} is a partial unification
in which quarks and leptons are placed in the same multiplet: the
gauge group $\SU(4)$ treats the three colours and the lepton number
as four ``colours,'' so a quark triplet plus a lepton singlet form a
$\mathbf{4}$. In the present context, the $\SU(4)_R$ flavour symmetry
of the $\Nf=4$ theory plays the role of the Pati--Salam colour
group. Gauging $\SU(3)_c \times \UU(1)_{B-L} \subset \SU(4)_R$ gives:
\begin{equation}
  \bar{\mathbf{4}} \;\to\; \bar{\mathbf{3}}_{-1/3}
    \,+\, \mathbf{1}_{+1}\,:
  \quad\text{3 quarks + 1 lepton}\,.
  \label{eq:patisalam}
\end{equation}
The baryon $B_i$ thus decomposes into three colour-triplet states
(identified with right-handed quarks) and one singlet (a
right-handed neutrino). The remaining SM gauge factors
$\SU(2)_L \times \UU(1)_Y$ are embedded in
$\SU(4)_L$ via $\mathbf{4}_L \to
\mathbf{2}_{-1/2} \oplus \mathbf{1}_{+1} \oplus \mathbf{1}_0$,
as in Section~\ref{sec:involution}.

The four baryoninos $\psi_{B_i}$ are the candidate right-handed
neutrinos. Under $\SU(4)\to\SU(3)\times\UU(1)$, the
$\bar{\mathbf{4}}$ decomposes into three $\nu_{R,i}$ plus one sterile
singlet. On the mesonic branch ($\langle B\rangle = \langle\tilde{B}\rangle = 0$),
the baryonino mass matrix is
$\mathcal{M}_{R}^{ij} = \partial^2 W/\partial B_i\partial\tilde{B}^j
= \tilde{M}^i{}_j/\Lambda^5$, where $\tilde{M}$ is the cofactor matrix
of~$M$.\footnote{The $R$-charge assignment $R(B)=R(\tilde B)=3/4$,
$R(M)=1/2$ selects $B\,M\,\tilde B$ as the marginal coupling.
The cofactor form $\tilde{M}/\Lambda^5$ arises after
canonically normalising the composite baryonic operators via
their K\"ahler metric, $K_B \sim |B|^2/\Lambda^4$.}
For diagonal $\langle M\rangle$, the $F$-term equations
$\prod_{k\ne i} v_k = m_i\,\Lambda^5$ imply
\begin{equation}
  \mathcal{M}_R^{ii} = m_i\,,
  \label{eq:baryonino_mass}
\end{equation}
so the baryonino masses equal the mass-deformation parameters. In particular,
$M_R(\nu_e) = m_e$ and $M_R(\nu_4) = m_t$. Meanwhile, no tree-level Dirac
coupling between mesinos and baryoninos exists on this branch
($\partial^2 W/\partial M\,\partial B \propto \langle\tilde{B}\rangle = 0$),
so neutrino masses arise only from radiative corrections, schematically
$m_\nu \sim m_e^2/(16\pi^2\,\Lambda) \approx 0.02\;\text{eV}$---the
correct order of magnitude. This provides a structural explanation for
the smallness of neutrino masses in the
sBootstrap~\cite{Rivero:2024epjc} (the ``SUSY-bootstrap'' framework
where hadronic SUSY relates mesons to fermions), without invoking a
separate seesaw scale.

A \emph{hybrid scenario} reconciles $\Nf=4$ with the $N=3$ uniqueness
of Section~\ref{sec:uniqueness}: at high energy, $\SU(3)$ SQCD with
4~flavours exhibits Pati--Salam structure. When the heaviest flavour
decouples ($m_4 = m_t \gg \Lambda$), holomorphic scale matching gives
$\Lambda_3^6 = \Lambda_4^5 \cdot m_t$, and the low-energy theory is
effective $\Nf = 3$ where the $\UU(3)$ charge structure applies.
The top quark plays the role of the spectator in the
turtle/elephant partition of~\cite{Rivero:2024epjc}: the one quark that
does not participate in the charge binding. With
$\Lambda_3 = f_\pi\sqrt{3}/2 \approx 80\MeV$
(Section~\ref{sec:lagrangian}) and $m_t = 173\GeV$, scale matching
gives $\Lambda_4 \approx 17\MeV$.
This scale is well below $\Lambda_{\rm QCD} \approx 220\MeV$ and
might seem dangerously low. However, the family $\SU(3)_F$ is a
\emph{separate} gauge group from QCD: its strong coupling scale
governs the composite Higgs/lepton masses, not hadronic binding.
The physical role of $\Lambda_4$ is to set the
normalisation of the $\Nf=4$ confined superpotential relative to the
mass deformations. Since $\Lambda_4 \ll m_t$, the $\Nf=4$ theory is
deeply in its confined phase when the fourth flavour is integrated out,
and the transition to the $\Nf=3$ effective theory is smooth.

\paragraph{PMNS mixing: an open problem.}
A significant phenomenological limitation of the mesonic-branch
construction is the neutrino mixing pattern. On this branch,
$\mathcal{M}_R$ is diagonal~\eqref{eq:baryonino_mass} and
the Dirac coupling vanishes at tree level, so
the PMNS matrix is the identity to all orders in perturbation theory
within the superpotential.
This is in clear conflict with the observed large lepton mixing
angles ($\theta_{12} \approx 34^\circ$, $\theta_{23} \approx 49^\circ$,
$\theta_{13} \approx 8.5^\circ$~\cite{NuFIT:2024}).
Potential sources of large PMNS mixing include:
(i)~K\"ahler-potential rotations, which are unconstrained in the
strongly-coupled regime and can be $O(1)$;
(ii)~SUSY-breaking contributions from the baryonic branch,
where $\langle B\rangle \neq 0$ generates off-diagonal Dirac couplings;
(iii)~threshold corrections at the $\Nf=4 \to \Nf=3$ matching scale.
We do not claim to explain PMNS mixing in this paper and leave its
computation as the most pressing open problem of the framework.


%======================================================================
\section{The Effective Lagrangian}
\label{sec:lagrangian}
%======================================================================

Below the family confinement scale~$\Lambda$, the sBootstrap effective
theory is specified by a single chiral superfield~$M^{\alpha}{}_i$
(the meson matrix), plus the elementary SM fermions. The UV theory has
gauge group $\SU(3)_F \times \SU(3)_c \times \SU(2)_L \times
\UU(1)_Y$, with preons $Q^\alpha_a \in (\mathbf{3}_F,\,
\mathbf{3}_L)$ and $\tilde{Q}_{a,i} \in (\bar{\mathbf{3}}_F,\,
\mathbf{1}_{\rm SM})$. After $\SU(3)_F$ confinement:

\paragraph{Superpotential.}
\begin{equation}
  W = \frac{\det M - B\tilde{B}}{\Lambda^3} + m_i\,M^i{}_i
    + \frac{y_\ell}{\Lambda}\,L_i\,e^c_j\,M^{1,2}{}_k\,\epsilon^{ijk}
    + \frac{y_u}{\Lambda^3}\,q_L\,u^c\,\adj(M)^{1,2}
    + \frac{y_d}{\Lambda}\,q_L\,d^c\,M^{1,2}\,,
  \label{eq:Wfull}
\end{equation}
where $M^{1,2}$ denotes the doublet ($Y{=}{-}1/2$) components and
$\adj(M)^{1,2}$ the doublet ($Y{=}{+}1/2$) components of the adjugate.
The lepton Yukawa is automatic: $L_i$ is the mesino
$\psi_{M^{1,2}{}_i}$ and $e^c_j$ is $\psi_{M^3{}_j}$, so the
$y_\ell$~term is contained in $\det M/\Lambda^3$ as a meson
self-coupling. (In the confined Hessian, the doublet--singlet block is
antisymmetric in generation indices; the symmetric SUSY-breaking
contributions that complete the physical mass matrix are discussed
below.) Quark Yukawas arise from partial compositeness mixing.

\paragraph{Field content.}
\begin{center}
\begin{tabular}{lcccl}
\toprule
Field & $\SU(3)_c$ & $\SU(2)_L$ & $\UU(1)_Y$ & Origin \\
\midrule
$L_i = \psi_{M^{1,2}{}_i}$ & $\mathbf{1}$ & $\mathbf{2}$ & $-1/2$ & mesino (composite) \\
$e^c_i = \psi_{M^3{}_i}$ & $\mathbf{1}$ & $\mathbf{1}$ & $+1$ & mesino (composite) \\
$\nu_R = \psi_B$ & $\mathbf{1}$ & $\mathbf{1}$ & $0$ & baryonino (composite) \\
$H_d = M^{1,2}$ & $\mathbf{1}$ & $\mathbf{2}$ & $-1/2$ & meson scalar \\
$H_u = \adj(M)^{1,2}$ & $\mathbf{1}$ & $\mathbf{2}$ & $+1/2$ & adjugate (function of $M$) \\
$q_L$, $u^c$, $d^c$ & $\mathbf{3}$, $\bar{\mathbf{3}}$, $\bar{\mathbf{3}}$ & --- & --- & elementary \\
\bottomrule
\end{tabular}
\end{center}

\paragraph{Key structural features.}
\begin{enumerate}
  \item \emph{One chiral field}: $\adj(M)$ is a quadratic polynomial in
    $M$; the superpotential~\eqref{eq:Wfull} is holomorphic in~$M$ alone.
    No independent $H_u$ superfield is needed.
  \item \emph{No free Yukawas}: The mass parameter $m_i$ simultaneously
    sets the UV mass deformation, the effective Yukawa coupling, and the
    meson VEV. The Yukawa hierarchy \emph{is} the family charge hierarchy.
  \item \emph{Determined $\tan\beta$}: At the superpotential level, the
    VEV ratio
    $\langle\adj(M)\rangle/\langle M\rangle = \Lambda^6/\langle M\rangle^2$
    is fixed by the quantum constraint. However, $\adj(M)$ is a
    composite operator, quadratic in~$M$; its physical VEV depends on
    the K\"ahler metric $K(M,M^\dagger)$, which is not protected by
    holomorphy. Whether $v_u/v_d$ remains fully determined at the
    quantum level requires computing the K\"ahler potential on the
    constraint surface---an open problem (cf.\ Table~\ref{tab:status}).
  \item \emph{Fermi scale}: The meson condensate
    $\langle M_\tau\rangle = \Lambda^2 m_\tau/P_\ell^{1/3}$ relates
    $v_{\rm EW}$ to $\Lambda$ through the K\"ahler normalization~$Z$
    (see~\S\ref{sec:discussion}). Matching
    $\Lambda = f_\pi\sqrt{3}/2 \approx 80\MeV$.
\end{enumerate}


%======================================================================
\section{Numerical Results}
\label{sec:numerics}
%======================================================================

\subsection{Charged leptons}

With $K = 2/3$ exact, the two lighter masses $m_e$ and~$m_\mu$
fix $k_{\rm lep} z_0^2 = 313.8\MeV$ and $\alpha_{\rm lep} =
0.22222(1)\;\text{rad}$; the third mass~$m_\tau$ is then a
\emph{prediction}.  The charge ansatz
\eqref{eq:massansatz}--\eqref{eq:zi} gives:

\begin{table}[ht]
\centering
\begin{tabular}{lccc}
\toprule
  & Predicted & PDG (MeV) & Deviation \\
\midrule
$m_e$ & 0.511 & 0.511 & input \\
$m_\mu$ & 105.66 & 105.66 & input \\
$m_\tau$ & 1776.97 & $1776.86 \pm 0.12$ & $+0.91\sigma$ \\
\bottomrule
\end{tabular}
\caption{Charged-lepton masses from the $\UU(3)$ charge ansatz.
The two inputs $(m_e, m_\mu)$ fix $k$ and~$\alpha$; the
predicted~$m_\tau$ deviates by $0.91\sigma$ from PDG.}
\label{tab:leptons}
\end{table}

\subsection{Quark sector and CKM elements}

\begin{table}[ht]
\centering
\begin{tabular}{lcccc}
\toprule
  Observable & Method & Predicted & PDG & Dev. \\
\midrule
$m_u + m_d$ & zero-charge ($\alpha{=}\pi/4$) & $6.7\MeV$ & $6.8\MeV$ & $1.4\%$ \\
$m_c$ & charge chain & $1301\MeV$ & $1270\MeV$ & $2.4\%$ \\
$m_b$ & inverse sector & $4567\MeV$ & $4180\MeV$ & $9.3\%$ \\
$m_t$ & standard chain & $180\GeV$ & $172.8\GeV$ & $3.9\%$ \\
$|V_{us}|$ & $\sqrt{m_d/m_s}$ & $0.2236$ & $0.2243 \pm 0.0008$ & $0.9\sigma$ \\
$|V_{cb}|$ & $1/(5R)$ & $0.0417$ & $0.0422 \pm 0.0008$ & $0.6\sigma$ \\
$|V_{ub}|$ & $2\sqrt{2}/(7R^3)$ & $0.00367$ & $0.00368 \pm 0.00011$ & $0.1\sigma$ \\
$\delta_{CP}$ & $\arctan\sqrt{R}$ & $65.4^\circ$ & $65.7 \pm 3.4^\circ$ & $0.1\sigma$ \\
$J$ & from above & $3.04{\times}10^{-5}$ & $3.08 \pm 0.15{\times}10^{-5}$ & $0.3\sigma$ \\
$\sin 2\beta$ & from above & $0.723$ & $0.699 \pm 0.017$ & $1.4\sigma$ \\
\bottomrule
\end{tabular}
\caption{Quark sector and CKM elements.
Inputs: $m_d = 4.67\MeV$, $m_s = 93.4\MeV$
($\overline{\rm MS}$ at~$2\GeV$). Here $R = (5+\sqrt{21})/2$.
The CKM expressions involving~$R$ are \emph{empirical pattern
identifications}; derivations from the operator basis remain open
(see text). The $m_b$ and~$m_t$ entries have elasticities
$\varepsilon \approx 6$ and~4 to~$m_d$; at $m_d = 4.60\MeV$ they
shift to $4180$ and~$170\GeV$ respectively.}
\label{tab:quarks}
\end{table}

\paragraph{Chain structure.}
The quark predictions in Table~\ref{tab:quarks} arise from a chain of
overlapping Koide triples~\cite{Rivero:2024epjc} that splits into two
regimes:

\noindent
\emph{Fermi side} (above the SUSY-breaking scale):
\begin{equation}
  (t,\,b,\,c):\quad K(m_t,\,m_b,\,m_c) \approx \tfrac{2}{3}
  \qquad\text{(standard Koide, $\overline{\rm MS}$ masses)}\,.
  \label{eq:chain_tbc}
\end{equation}

\noindent
\emph{Chiral side} (below the breaking scale):
\begin{equation}
  (c,\,s,\,0):\quad \frac{m_c}{m_s} = (2{+}\sqrt{3})^2\,;
  \qquad
  (s,\,0,\,m_u{+}m_d):\quad \frac{m_s}{m_u{+}m_d} = (2{+}\sqrt{3})^2\,.
  \label{eq:chain_chiral}
\end{equation}
The chiral-side links use zero-charge triples ($z_2 = 0$, \ie,
$\alpha = \pi/4$), where one Koide charge vanishes and
$K = 2/3$ reduces to a mass ratio $(2{+}\sqrt{3})^2 = 13.93$.
The lightest quark entry uses the \emph{isospin
reparametrization} $u \to 0$, $d \to m_u + m_d$: only the
total light-quark mass enters, because below $\Lambda_{\rm QCD}$
the individual $m_u$ and $m_d$ are unobservable---the
Gell-Mann--Oakes--Renner relation
$m_\pi^2 = B_0(m_u{+}m_d)$~\cite{GellMannOakesRenner:1968}
determines only the sum.
This reparametrization is natural in the sBootstrap interpretation,
where Koide is a scalar (meson) mass formula:
$(c,s,0)$ connects to the zero-charge meson Koide
$K(0,\, m_\pi,\, m_D) = 2/3$ (with $\sqrt{m_D/m_\pi} \approx
2{+}\sqrt{3}$ to~$2\%$).
The two sides are connected by the tuple $(b,c,s)$ with a
\emph{signed} square root: $K(\sqrt{m_b} + \sqrt{m_c} - \sqrt{m_s})
= 0.675$, deviating from~$2/3$ by~$1.2\%$.
The negative sign on $\sqrt{m_s}$ corresponds to
$z_0 + z_s < 0$ in the charge ansatz---the sector angle has crossed
the critical value where the lightest charge flips sign.
The complete chain is therefore
\begin{equation}
  (t,b,c) \;-\; (b,c,s)^{\pm} \;-\; (c,s,0) \;-\; (s,0,m_u{+}m_d)\,,
  \label{eq:fullchain}
\end{equation}
with $(b,c,s)$ as the bridge link. The signed tuple
hits $K = 2/3$ exactly at $m_s \approx 84\MeV$---between
the PDG value $m_s(2\GeV) = 93.4\MeV$ and the higher-scale
estimate $m_s(m_b) \approx 81\MeV$---suggesting that its
natural evaluation scale is $\mu \sim 3$--$5\GeV$.
At~$1.2\%$ the bridge tuple is not compelling on its own,
but its structural role as the link between Fermi and chiral
regimes is significant.

The Fermi/chiral split coincides with the transition from
fermionic to scalar mass formulae: above the breaking scale
the SUSY mass relation $m_{\rm fermion} \approx m_{\rm scalar}$
makes the chain applicable to quark masses directly; below it,
the chain applies to meson masses, with the isospin
reparametrization reflecting the fact that chiral symmetry
breaking groups $u$ and~$d$ into a single Goldstone sector.

\paragraph{IR/UV scale preference.}
The two Yukawa operators of Section~\ref{sec:involution}---$M$
(meson VEV, masses $\propto 1/m_{\rm UV}$) and $\adj(M)$
(adjugate, masses $\propto m_{\rm UV}$)---make a structural
prediction about scale dependence: the $W_1$~sector (species
coupling through~$M$) should be evaluated at \emph{infrared}
scales (pole masses or low $\overline{\rm MS}$ scales), while
the $W_2$~sector (species coupling through $\adj(M)$) should be
evaluated at \emph{ultraviolet} scales (high $\overline{\rm MS}$).
Using the AHS running masses~\cite{AntuschHinzeSaad:2025}:
\begin{center}
\begin{tabular}{lccl}
\toprule
Tuple & Best $K - 2/3$ & At $M_Z$ & Preferred scale \\
\midrule
$K(e,\mu,\tau)$ & $0.001\%$ (pole) & $0.18\%$ & IR \\
$K^{-1}(d,s,b)$ & $0.25\%$ ($m_i$) & $0.13\%$ & UV \\
$K(t,b,c)$ & $0.4\%$ ($m_i$) & $8\%$ & IR \\
$K(c,s,0)$ & $0.3\%$ ($m_i$) & $2.6\%$ & IR \\
\bottomrule
\end{tabular}
\end{center}
Every chain tuple except the inverse down quarks worsens at $M_Z$,
confirming they belong to the $W_1$~sector ($M \propto 1/m$).
The inverse Koide \emph{improves} at $M_Z$ and passes through
$2/3$ near~$1\TeV$ (\S\ref{sec:running}), consistent with
$\adj(M) \propto m$ evaluated at UV scales. This IR/UV
dichotomy is a structural consequence of the two operators,
not a tuning of evaluation scales.

The empirical relation $|V_{cb}| = 1/(R^2 + 1) = (5-\sqrt{21})/10$
is algebraically natural given $R^2 + 1 = 5R$ (a consequence of
$R^2 - 5R + 1 = 0$), but a derivation from the operator basis of
Section~\ref{sec:nf4} has not been established.
Combined with $|V_{us}| = \sqrt{m_d/m_s}$, this gives the Wolfenstein
$A$~parameter: $A = |V_{cb}|/|V_{us}|^2 = m_s/(5R\,m_d) = 0.835$,
compared to the PDG value $A = 0.836 \pm 0.015$ ($0.1\%$).

The CP phase $\delta_{CP} = \arctan\sqrt{R}$ is equivalent to
$\cos^2 2\delta = 3/7$, or $\sin 2\delta = 2/\sqrt{7}$, since
$(R{+}1)^2 = 7R$. This is not an independent prediction but a
consequence of the $N=3$ uniqueness:
since $(R{-}1)^2 = NR$ and $(R{+}1)^2 = (N{+}4)R$ (from the self-dual
quadratic), one has $\cos^2 2\delta = N/(N{+}4)$. This equals
$N/(2N{+}1)$ only when $N{+}4 = 2N{+}1$, \ie, $N=3$---the same
condition as $D_{\rm alg} = D_{\rm dyn}$.
The Jarlskog invariant has the closed form
$J \approx 2\sqrt{2}/(5 \times 7^{5/4})\,\sqrt{m_d(m_s{-}m_d)}/m_s
\times R^{-15/4} = 3.04 \times 10^{-5}$ (PDG: $3.08 \times 10^{-5}$).
Figure~\ref{fig:ut} shows the predicted unitarity triangle, with the
apex at $(\bar\rho, \bar\eta) = (0.164, 0.358)$ and the three angles
$\alpha = 92.0^\circ$, $\beta = 22.6^\circ$, $\gamma = 65.4^\circ$,
all within $1.5\sigma$ of PDG constraints.

\begin{figure}[ht]
\centering
\includegraphics[width=0.7\textwidth]{figures/ut_triangle.pdf}
\caption{The CKM unitarity triangle from the framework's four inputs.
The shaded bands show PDG constraints on $\sin 2\beta$ and~$\gamma$.
The predicted apex (black dot) lies at the intersection of the
$R_b$~and $R_t$ circles.}
\label{fig:ut}
\end{figure}

\subsection{Sensitivity and parameter counting}
\label{sec:paramcount}

\paragraph{Input--output accounting.}
The framework has four continuous input parameters:
\begin{enumerate}
  \item $k$ (overall mass scale) --- fixed by $m_e$;
  \item $\alpha$ (sector angle) --- fixed by $m_\mu$;
  \item $m_d$ ($\overline{\rm MS}$ at $2\GeV$) --- taken from PDG;
  \item $m_s$ ($\overline{\rm MS}$ at $2\GeV$) --- taken from PDG.
\end{enumerate}
In addition, the framework relies on discrete structural choices:
the charge ansatz $m_i = k(z_0+z_i)^2$, the partial compositeness
exponents $\gamma_\ell = 0$ and $\gamma_d = 1$, and the
K\"ahler normalization~$Z$ required to match $v_{\rm EW}$.
These are not continuous parameters but foundational assumptions
(Table~\ref{tab:status}).

From these four inputs, the framework produces the following
genuine predictions (outputs that are not used to fix any input):
\begin{itemize}
  \item $m_\tau$ (from $k$, $\alpha$, and $K=2/3$): $+0.91\sigma$.
  \item $m_c = m_s(2{+}\sqrt{3})^2$: $2.4\%$ above PDG.
  \item $m_b$ from $K^{-1}(d,s,b) = 2/3$: sensitive to $m_d$
    ($\varepsilon = 6.3$).
  \item $m_t$ from the chain: sensitive to $m_d$
    ($\varepsilon \approx 4$).
  \item $m_u + m_d$ from $K(0, s, m_u{+}m_d) = 2/3$: $1.4\%$.
  \item $|V_{us}| = \sqrt{m_d/m_s}$: $0.9\sigma$.
\end{itemize}
The CKM elements $|V_{cb}|$, $|V_{ub}|$, and $\delta_{CP}$
depend only on the self-dual ratio~$R = (5+\sqrt{21})/2$ and are
\emph{independent of all four continuous inputs}. However, these
are currently \emph{empirical pattern identifications}, not
derivations from the operator basis. The Jarlskog invariant~$J$
and $\sin 2\beta$ are derived quantities (not independent of
$|V_{cb}|$, $|V_{ub}|$, $\delta_{CP}$).

Figure~\ref{fig:sensitivity} shows the sensitivity of $m_b$
and $m_t$ to the input $m_d$.

The inverse Koide formula amplifies the $m_d$ uncertainty by a factor
of~6 in~$m_b$ and~4 in~$m_t$ (Fig.~\ref{fig:sensitivity}); given the
PDG range $m_d = 4.7 \pm 0.4\MeV$, the predicted~$m_b$ spans
${\sim}\,3000$--$7000\MeV$. Setting $m_d = 4.60\MeV$ (within
$0.24\sigma$ of PDG central) simultaneously gives
$m_b = 4180\MeV$ and $m_t = 170\GeV$ ($-1.7\%$); the quark
predictions are consistent with experiment at every $m_d$ in
the PDG range.

\begin{figure}[ht]
\centering
\includegraphics[width=0.55\textwidth]{figures/sensitivity_plot.pdf}
\caption{Inverse Koide predictions for $m_b$ (top) and the chain
prediction for~$m_t$ (bottom) as functions of the input~$m_d$.
The red bands show PDG central values $\pm 1\sigma$;
the grey band is the PDG range for~$m_d$.
The red dot marks $m_d = 4.60\MeV$, where $m_b$ and~$m_t$
simultaneously match PDG.}
\label{fig:sensitivity}
\end{figure}

In contrast, $m_c = m_s\,(2{+}\sqrt{3})^2$
has unit elasticity to~$m_s$ and is the most predictive quark relation.
At PDG central $m_s = 93.4\MeV$, the chain gives
$m_c = 1301\MeV$ ($1.5\sigma$ above PDG); at the Rivero value
$m_s = 95\MeV$, it gives $m_c = 1323\MeV$ ($2.7\sigma$). This is
the largest tension in the framework, and it sharpens considerably
once scale dependence is examined carefully.

\paragraph{Holomorphic masses versus physical masses.}
The $\overline{\rm MS}$ anomalous dimension~$\gamma_m$ is
flavour-independent at all known loop orders (verified through
five loops~\cite{Baikov:2014qja}), so the ratio
$m_c(\mu)/m_s(\mu)$ \emph{does not run} in QCD\@.
The most precise lattice determination gives
$m_c/m_s = 11.783 \pm 0.025$~\cite{FNAL:2018mc_ms}, which is
$18\%$ below $(2{+}\sqrt{3})^2 = 13.93$.
The frequently quoted ``agreement'' at the level of $13.6$ arises
only from comparing $m_s$ at~$2\GeV$ with $m_c$ at~$m_c$---an
inconsistent mixed-scale comparison that accidentally compensates for
the discrepancy.

This tension has a natural interpretation in the present framework.
The mass deformations~$m_i$ entering the superpotential
$W_{\rm tree} = \Tr(m\,M)$ are \emph{holomorphic parameters}:
they are protected by the non-renormalisation theorem and do not run
at any scale~\cite{Intriligator:1995au,
ArkaniHamedMurayama:1997}.  Physical (pole or $\overline{\rm MS}$)
masses, by contrast, are related to the holomorphic ones through the
K\"ahler metric:
\begin{equation}
  m_i^{\rm phys} \;=\; \frac{m_i^{\rm hol}}{\sqrt{Z_i}}\,,
  \label{eq:kahler_matching}
\end{equation}
where $Z_i$ is the wavefunction normalisation of the $i$-th
composite.  In a strongly coupled s-confining theory, the~$Z_i$
are generation-dependent---different mass deformations break
$\SU(3)_F$ differently---and non-perturbatively unknown.

If the chain relation $(2{+}\sqrt{3})^2$ is a statement about
holomorphic mass ratios, then the $18\%$ discrepancy with the
physical $m_c/m_s$ requires $Z_c/Z_s \approx 1.37$: a $37\%$
difference in wavefunction normalisation between charm and strange
composites.  For strongly coupled scalars this is not unreasonable,
and it connects the chain tension directly to the central open
problem of computing the K\"ahler metric (Section~\ref{sec:discussion}).

By contrast, the \emph{lepton} Koide relation compares pole masses
of a single generation-universal composite type (mesinos), where
gauge corrections to~$Z_i$ are generation-blind at leading order
and the holomorphic prediction transfers to physical masses with
only $\mathcal{O}(\alpha/4\pi)$ corrections---consistent with
the observed $0.04\%$ accuracy.

(The 2024 PDG uncertainties on $m_b$ and~$m_s$ are roughly
10~times smaller than the 2022 values, making the inverse Koide
deviation of~$0.25\%$ a ${\sim}\,1\sigma$ tension at the conventional
scales.  This tension is relieved at $M_Z$, where the deviation
falls to~$0.13\%$.)
Figure~\ref{fig:pulls} summarises the pulls of all CKM predictions
and the lepton deviations.

\begin{figure}[ht]
\centering
\includegraphics[width=0.85\textwidth]{figures/pull_plot.pdf}
\caption{Left: pulls (prediction~$-$~PDG central)/$\sigma_{\rm exp}$
for the CKM observables. Green bars are
parameter-free predictions (from~$R$ only); orange bars depend on the
quark mass inputs $m_d$,~$m_s$. Right: the predicted $m_\tau$ ($+0.9\sigma$),
with $m_e$ and~$m_\mu$ used as inputs to fix $k$ and~$\alpha$.}
\label{fig:pulls}
\end{figure}

\subsection{The pion decay constant and the confinement scale}

The family confinement scale is identified as
$\Lambda = f_\pi\sqrt{3}/2 \approx 80\MeV$, where
$f_\pi = 92.07\MeV$. The quantum constraint
$\det M = \Lambda^6$ implies
$(m_e m_\mu m_\tau)^{1/3} = y_1 \Lambda$ with
$y_1 = 1/\sqrt{3}$; combined with the confinement-scale
identification, this fixes
\begin{equation}
  f_\pi = 2(m_e m_\mu m_\tau)^{1/3}
  \label{eq:fpi_P}
\end{equation}
to~$0.56\%$---a relation between the pion decay constant and the
lepton mass product that is a consequence of the framework, not an
additional input.

\subsection{Lepton--quark angle shift}

The lepton sector angle $\alpha_{\rm lep} =
0.22222(1)\;\text{rad} = 12.73^\circ$ and the inverse down-quark sector
angle $\alpha_{\rm inv} = 10.59^\circ$ (from $K^{-1}(d,s,b) = 2/3$
at the physical masses) differ by $\Delta\alpha = \alpha_{\rm lep} -
\alpha_{\rm inv} = 2.14^\circ$. Empirically,
\begin{equation}
  \tan(\Delta\alpha) \;=\; \frac{\pi}{84}
  \;=\; \frac{\pi}{4\,D}
  \label{eq:angleshift}
\end{equation}
to~$1.6\%$, where $D = 21$ is the discriminant.  If exact,
this formula determines the quark mass ratio
$m_d/m_s = 0.0492$ (PDG: $0.050 \pm 0.003$), reducing the parameter
count from four to three: only $(k, \alpha, m_s)$ are needed.
No derivation of \eqref{eq:angleshift} is known.

\subsection{Charge-plane geometry}
\label{sec:chargeplane}

The Koide parametrisation $\sqrt{m_i} = \sqrt{k}\,(z_0 + z_i)$ with
$\sum z_i = 0$ places the charges $z_i$ in a two-dimensional plane,
whose orientation is specified by the sector angle~$\alpha$.
The charge-plane angles (mod~$120^\circ$) of the SM Koide triples
cluster within ${\sim}\,5.5^\circ$ of the zero-mass $K = 2/3$
fixed point at $\alpha = 15^\circ$ (corresponding to
the mass ratio $r_0 = (2+\sqrt{3})^2 = 13.93$).
This clustering is a consequence of all SM Koide tuples having
their dominant mass ratio close to~$r_0$.

\subsection{Scale dependence}
\label{sec:running}

Using the running Yukawa couplings of Antusch, Hinze, and
Saad~\cite{AntuschHinzeSaad:2025} (SM, $\overline{\rm MS}$~scheme,
2024 PDG inputs), the Koide ratio for the charged leptons is
$K_{\rm lep} = 0.6678$ at all scales from $M_Z$ to $10^7\GeV$
($0.17\%$ from $2/3$). The larger deviation compared to the pole-mass
value ($0.001\%$) reflects QED running corrections
$\delta m \propto \alpha\,\log(\mu^2/m_{\rm pole}^2)$, which affect
each lepton differently. The approximate scale invariance of~$K$ under
RG flow supports Sumino's
observation~\cite{Sumino:2009} that gauge interactions
approximately preserve the Koide relation.

For the inverse Koide ratio of the down quarks, the running reveals an
interesting structure: $K(1/m_d,\,1/m_s,\,1/m_b)$ passes through
$2/3$ at $\mu \approx 1\TeV$, with the central value
$K - 2/3 = +2\times 10^{-6}$ at $10^3\GeV$. The deviation grows to
$+0.13\%$ at $M_Z$ and $-0.2\%$ at $10^7\GeV$. Although the
current experimental uncertainties on $m_d$ and $m_s$ imply a
${\sim}\,0.4\%$ band on~$K$, the trend is robust: the inverse Koide
for down quarks is best satisfied near the TeV scale, potentially
connected to the effective adjoint mass scale~$\mu_\Phi$ of the
emergent doublet structure (\S\ref{sec:discussion}).
Figure~\ref{fig:running} shows both Koide ratios as a function of
scale, using the AHS running Yukawa couplings in the SM
$\overline{\rm MS}$~scheme.

\begin{figure}[ht]
\centering
\includegraphics[width=0.7\textwidth]{figures/running_koide.pdf}
\caption{Koide ratios vs.\ renormalization scale in the SM.
Blue: $K(e,\mu,\tau)$; red: $K^{-1}(d,s,b)$; dashed:
$K=2/3$. The shaded band shows the uncertainty from $m_d$.
Data from Antusch, Hinze, and Saad~\cite{AntuschHinzeSaad:2025}.}
\label{fig:running}
\end{figure}

%======================================================================
\section{Discussion}
\label{sec:discussion}
%======================================================================

\subsection{What is derived}

The charge ratio $K = 2/3$ follows from $\UU(3)$ normalization. The
standard--inverse duality follows from the $F$-term structure of the
s-confined theory: the adjugate $\adj(M) \propto m$ and the meson
$M \propto 1/m$ are the two natural Yukawa operators at the
confinement scale, leading to standard and inverse Koide for the
respective SM species. The duality is identified with the
$\mathbf{15} \leftrightarrow \overline{\mathbf{15}}$ exchange of
$\SU(5)$ and the Galois conjugation of $\mathbb{Q}(\sqrt{21})$.
Three generations is uniquely selected by discriminant matching. These
are exact results contingent on the charge ansatz
\eqref{eq:massansatz}---itself uniquely selected by the
direction-independence of $K$ (Section~\ref{sec:u3})---and the family
symmetry structure.

\subsection{The emergent doublet structure and the critical scale}

The confined spectrum contains two composite operators per
family---$M_i$ and $\adj(M)_i$---whose mass eigenvalues have
opposite charge dependence (direct and inverse branches).
This pairing has the mathematical form of an $\mathcal{N}=2$
hypermultiplet~\cite{SeibergWitten:1994a,SeibergWitten:1994b},
but is emergent from confinement plus isospin, not a UV input.
The coupling between the two branches can be parametrised by
an effective quartic meson interaction with scale~$\mu_\Phi$:
\begin{equation}
  W_{\rm eff} \;\supset\; -\frac{1}{2\mu_\Phi}\Tr(M^2)\,,
  \qquad M = Q\tilde Q\,.
  \label{eq:quartic}
\end{equation}
The full superpotential in the s-confined description becomes
\begin{equation}
  W = \frac{\det M - B\tilde B}{\Lambda^3}
    + m_i M^i{}_i
    - \frac{1}{2\mu_\Phi}\Tr(M^2)\,,
  \label{eq:N2flow}
\end{equation}
with modified $F$-term equation
\begin{equation}
  \frac{\adj(M)_i}{\Lambda^3} + m_i - \frac{M_i}{\mu_\Phi} = 0\,.
  \label{eq:Fterm_mu}
\end{equation}
The limit $\mu_\Phi\to\infty$ recovers the pure $\mathcal{N}=1$ theory.
At finite~$\mu_\Phi$, the Koide vacuum exists only above a critical
scale
\begin{equation}
  \mu_{\Phi,\rm crit} = \frac{4\,\Lambda^2\,P_\ell^{1/3}}{m_e^2}
  \approx 4.5\TeV\,.
  \label{eq:mucrit}
\end{equation}
Below this scale, the $\Tr(M^2)$ correction overwhelms
the quantum constraint for the electron meson VEV.
The true phase boundary lies at
$\mu_{\Phi,\rm bif} = \mu_{\Phi,\rm crit}/2^{1/3} \approx 3.5\TeV$.
Above $\mu_{\Phi,\rm bif}$, the Koide ratio converges as
$K - 2/3 \approx A/\mu_\Phi$ where $A/\mu_{\Phi,\rm crit} = (z_0^2 - z_e^2)/6$;
the observed accuracy $K = 2/3$ to $0.001\%$ requires $\mu_\Phi \gtrsim 37\TeV$.
This constrains the effective inter-branch coupling
scale in the family sector to be at least $\sim\!40\TeV$.
Independently, the inverse down-quark Koide ratio
$K^{-1}(d,s,b)$ crosses $2/3$ at $\mu \approx 1$--$1.6\TeV$
under standard $\overline{\rm MS}$ running~\cite{AntuschHinzeSaad:2025},
the same order of magnitude as~$\mu_{\rm crit}$.

\subsection{The Fermi scale and K\"ahler normalization}

The meson condensate $\langle M_i\rangle$ has mass dimension~2.
The canonical Higgs field is
$h_i = Z^{1/2}\,M_i$ where $Z = K_{M\bar M}$ is the K\"ahler
wave-function renormalization:
\begin{equation}
  v_{\rm EW}^2 = Z\,\langle M_\tau\rangle^2
  = Z\,\frac{\Lambda^4\,m_\tau^2}{P_\ell^{2/3}}\,.
  \label{eq:vew}
\end{equation}
The three doublet VEVs are hierarchical:
$v_i/v_\tau = m_{\ell,i}/m_\tau$, so
$v_\tau/\sqrt{\sum v_i^2} = 0.998$---the electroweak VEV is
dominated by the heaviest lepton.
Matching $v_{\rm EW} = 246.2\GeV$ determines $Z$:
\begin{equation}
  Z^{1/2} = \frac{v_{\rm EW}\,P_\ell^{1/3}}{\Lambda^2\,m_\tau}\,.
  \label{eq:Zfactor}
\end{equation}
If $\Lambda = f_\pi\sqrt{3}/2 \approx 79.7\MeV$, the canonical
K\"ahler normalization $Z = 1/\Lambda^2$ gives
$v_{\rm EW} = 3.09\GeV$---two orders of magnitude below the
observed value. The discrepancy requires a large anomalous K\"ahler
correction $Z \gg 1/\Lambda^2$, characteristic of the
strongly-coupled confined phase.
By na\"{\i}ve dimensional analysis (NDA), the K\"ahler metric for a
composite of dimension~$d$ scales as
$Z \sim (4\pi)^{2(d-1)}/\Lambda^{2(d-1)}$~\cite{Luty:1998}; for
$d=2$ this gives $Z \sim (4\pi)^2/\Lambda^2$, an enhancement
of $\sim 160$ over canonical.
Both estimates are parametrically consistent with the required
$Z/Z_{\rm canon} \approx 6{,}300$, but off by an $O(1)$ coefficient that
cannot be fixed without a non-perturbative computation---either from
the Seiberg--Witten prepotential on the Coulomb branch
(which is known for $\SU(3)$ with 3~flavours~\cite{Argyres:1995wt})
or from lattice SQCD.
Computing the K\"ahler potential on the
quantum-modified moduli space remains the most important
quantitative open problem in the framework.

\subsection{SUSY breaking from electroweak symmetry breaking}

The SUSY-preserving vacuum is diagonal:
$\langle M^\alpha{}_i\rangle = v_\alpha\,\delta_{\alpha i}$, which does
not break $\SU(2)_L$ and has $V = 0$.
When $\SU(2)_L$ is broken by giving the Higgs doublet components
VEVs in all three generations, the $F$-term system becomes
overconstrained: 18~real equations for 12~real variables.
SUSY necessarily breaks, with the breaking scale set by the
lepton masses:
\begin{equation}
  |F_{21}| = m_e\,, \qquad
  |F_{22}| = 0\,, \qquad
  |F_{23}| = m_\tau\,,
  \label{eq:Fterms_lepton}
\end{equation}
and $|F| = \sqrt{\sum |F_{2j}|^2} \approx m_\tau \approx 1.8\GeV$.
The goldstino is predominantly the $\tau$-mesino (since
$m_\tau/m_e \approx 3500$), and
$m_{3/2} = |F|/(\sqrt{3}\,M_{\rm Pl})
\approx 4\times 10^{-16}\GeV$.

The three Higgs doublets $H_i = M^{1,2}{}_i$ acquire hierarchical
VEVs $v_i \propto m_{\ell,i}$ from the adjugate mechanism: $v_\tau$
dominates ($v_\tau/\sqrt{\sum v_i^2} = 0.998$), so the SM Higgs is
predominantly the $\tau$-generation doublet.

\subsection{The scalar spectrum and the mild SUSY-breaking problem}
\label{sec:mild}

Standard SUSY-breaking constructions (gravity or gauge mediation)
generate soft scalar masses $m_{\rm soft} \sim F/M_{\rm mess}$ that
send superpartners to the TeV scale---orders of magnitude above their
Standard Model counterparts.  In the present framework, the situation
is structurally different.  The SUSY-breaking $F$-terms \emph{are}
the lepton masses ($|F| \approx m_\tau \approx 1.8\GeV$, not
$(100\TeV)^2$), and the ``superpartners'' of the leptons are the
composite Higgs doublets that live in the same chiral superfield~$M$.
The scalar--fermion splitting is controlled by~$F$ and the family
charges, not by a distant hidden sector.

\paragraph{Charge-aligned K\"ahler deformation.}
The holomorphic mass formula $m_i = k(z_0{+}z_i)^2$ is protected by
the non-renormalization theorem: any deformation that preserves the
superpotential preserves it exactly.  The leading SUSY-breaking
correction therefore enters through the K\"ahler potential.
The simplest charge-aligned operator is
\begin{equation}
  \Delta K \;=\; \frac{c}{M_*^2}\,\sum_i (z_0{+}z_i)^2\;
  |F_X|^2\;|\Phi_i|^2\,,
  \label{eq:DK_spurion}
\end{equation}
where $X$ is a SUSY-breaking spurion with $\langle F_X \rangle \neq 0$
and $M_*$ is the mediation scale.  Because the weight is
$(z_0{+}z_i)^2$---the same charge combination that determines~$m_i$---the
resulting soft scalar mass
\begin{equation}
  m^2_{\rm soft,\,i} \;=\; \varepsilon\,m_i\,,
  \qquad
  \varepsilon \equiv \frac{c\,|F_X|^2}{k\,M_*^2}
  \label{eq:msoft}
\end{equation}
is automatically aligned with the mass hierarchy.  The scalar-to-fermion
mass ratio is $\sqrt{1 + \varepsilon/m_i}$: for
$\varepsilon \ll m_e$, \emph{all} partners remain nearby.  The fermion
masses (and therefore the Koide formula) are unaffected at this order.

\paragraph{Self-mediated estimate.}
In the most economical scenario, the spurion is the meson sector itself
($|F_X| \approx m_\tau$, $M_* \approx \Lambda$):
\begin{equation}
  \varepsilon \;\sim\; \frac{c\,m_\tau^2}{k\,\Lambda^2}
  \;\approx\; c \times 0.26\MeV\,.
  \label{eq:eps_self}
\end{equation}
For a loop-suppressed coefficient $c \sim \alpha/4\pi \sim 10^{-2}$,
$\varepsilon \sim 3\keV$, giving sub-percent splittings for all
three generations---consistent with the observed precision of~$K$.
For~$c \sim O(1)$, $\varepsilon \sim m_e/2$: the $\tau$ and $\mu$
partners are split by~$<0.2\%$, while the electron partner receives an
$O(1)$ relative correction.  In either case the splitting is
\emph{mild}: no partner is pushed to the TeV scale.

\paragraph{Wavefunction mechanism.}
An equivalent viewpoint replaces the spurion with a family-dependent
wave-function renormalization $Z_i = Z\bigl[1 + \varepsilon_1
(z_0{+}z_i)^2\bigr]$, giving
$m_i^{\rm phys} = m_i^{\rm hol}/\sqrt{Z_i}
\approx \tilde m_i\bigl[1 - (\varepsilon_1/2k)\,\tilde m_i\bigr]$
at first order.  The correction is
proportional to~$m_i^2$ (quartic in charge) and preserves $K = 2/3$
to~$O(\varepsilon_1^2)$.  This is precisely the partial-compositeness
mechanism of the next subsection: the anomalous dimension~$\gamma_*$
controls~$\varepsilon_1$, with $\gamma_* = 0$ (universal~$Z$) for
leptons and $\gamma_* = 1$ for down quarks.

\paragraph{The supertrace constraint.}
For $F$-term breaking with canonical K\"ahler,
$\mathrm{STr}(\mathcal{M}^2) = \sum m^2_{\rm boson} - \sum m^2_{\rm fermion} = 0$.
This means the sum of scalar mass-squared corrections vanishes:
$\sum_i m^2_{\rm soft,\,i} = \varepsilon \sum_i m_i = 0$
cannot hold (since $m_i > 0$ and $\varepsilon > 0$), signalling that
the canonical K\"ahler is modified at $O(\varepsilon)$.  The leading
non-canonical correction is the family-dependent~$Z_i$ above, which
redistributes the supertrace across generations without violating the
sum rule.

\paragraph{Contrast with the MSSM.}
In the MSSM, the $\mu$-problem and the hierarchy problem arise because
soft masses must be tuned against $|\mu|^2$ to get $v_{\rm EW} \ll
M_{\rm Pl}$.  Here, both~$v_{\rm EW}$ and the soft masses are
determined by the \emph{same} family charges (eq.~\ref{eq:massansatz}):
the electroweak VEV is $v_{\rm EW} = \Lambda^2 m_\tau/P_\ell^{1/3}$
(eq.~\ref{eq:vew}), the soft masses are $\varepsilon \sim m_\tau^2/
(k\Lambda^2)$, and the Higgs mass arises from Coleman--Weinberg
corrections at $M_S \approx \mu_{\rm crit}$ (eq.~\ref{eq:higgs_CW}).
There is no separate hidden sector; all scales are tied to the lepton
mass product~$P_\ell$ and the confinement scale~$\Lambda$.

\paragraph{What remains open.}
Computing the coefficient~$c$ in~\eqref{eq:DK_spurion} from first
principles requires the K\"ahler metric on the quantum-modified moduli
space of the s-confined theory.  For the genus-0 Seiberg--Witten curve
relevant to~$\Nf = \Nc = 3$, this metric is not constrained by
holomorphy and receives non-perturbative corrections from M2-brane and
M5-brane instantons (Appendix~\ref{app:Mtheory}).  The full scalar
spectrum---three Higgs doublets, three charged sleptons, and the
baryonic modulus---depends additionally on the $\SU(2)_L \times
\UU(1)_Y$ $D$-terms, which enforce vacuum alignment but whose
quantitative contribution to the mass splittings is model-dependent.
Identifying a mildly broken SUSY sector that keeps superpartners
nearby---rather than sending them to the TeV scale as in standard
constructions---is the most pressing open problem of the framework.

\subsection{Partial compositeness and quark masses}

Quarks are elementary fields, external to the confining sector.
They acquire mass through \emph{partial
compositeness}~\cite{Kaplan:1991dc}: each elementary quark mixes
linearly with a composite operator of the same quantum numbers:
$\mathcal{L} \supset \lambda_L\,\bar q_L\,\mathcal{O}_R
+ \lambda_R\,\bar q_R\,\mathcal{O}_L + \text{h.c.}$
The mixing coefficients $\lambda_i$ are generation-dependent,
determined by anomalous dimensions at the confining fixed
point~\cite{Strassler:1995,NelsonStrassler:2000}.
The superpotential is holomorphic in $M$ and $\adj(M)$,
so the most general Yukawa operator is
$W \supset q\,\bar{q}\,M^a\,\adj(M)^b/\Lambda^n$ with
$a, b \in \mathbb{Z}_{\ge 0}$. On the constraint surface
$\adj(M) = \Lambda^6 M^{-1}$, the generation-dependent mass
scaling is $m_q \propto m_{\rm UV}^{a-b}$: the exponent
$p = a - b$ is automatically an integer. EFT power counting
selects the two leading operators: $(a,b) = (1,0)$ (renormalisable,
$p{=}{+}1$) and $(a,b) = (0,1)$ ($p{=}{-}1$), related by the
$\mathbf{15} \leftrightarrow \overline{\mathbf{15}}$ exchange.
The \emph{assignment} (why down quarks couple to $M$ rather than
$\adj(M)$) is fixed by hypercharge
($Y(M^{1,2}) = -\tfrac{1}{2}$, $Y(\adj(M)^{1,2}) = +\tfrac{1}{2}$).

\subsection{What remains open}
\label{sec:open}

\begin{enumerate}
  \item \textbf{The dynamical origin of~$\alpha_{\rm lep}$.}
    All $\SU(3)$-invariant potentials for the traceless diagonal field
    depend on~$\alpha$ only through $\cos 3\alpha$, whose extrema lie at
    $\alpha = 0$ and $\pi/3$.  The one-loop Coleman--Weinberg effective
    potential remains $\mathbb{Z}_3$-invariant
    and monotonic in~$\alpha$ at every renormalization scale.
    Fixing~$\alpha_{\rm lep}$
    therefore requires non-perturbative K\"ahler effects,
    cross-sector constraints, or input from a UV completion.

  \item \textbf{$|V_{ub}|$ and the exclusive/inclusive tension.}
    $|V_{ub}| = 2\sqrt{2}/((2N{+}1)R^3)$ gives $0.003674$,
    matching the \emph{exclusive} determination
    $|V_{ub}|_{\rm excl} = 0.003682 \pm 0.000110$ ($0.08\sigma$) but
    deviating from the inclusive value
    $|V_{ub}|_{\rm incl} = 0.00425 \pm 0.00010$ ($5.8\sigma$).
    If the exclusive/inclusive tension
    is resolved in favour of the inclusive value, the framework is
    falsified.

  \item \textbf{PMNS mixing.}
    The baryonino Majorana matrix $\mathcal{M}_R$ is
    diagonal on the mesonic branch, and the
    mesino--baryonino Dirac coupling vanishes at tree level.
    The most promising path is non-diagonal K\"ahler corrections
    to the baryonic sector: the composite operator $B = \epsilon\,QQQ$
    involves all three flavours symmetrically, so its K\"ahler metric
    need not inherit the meson hierarchy.  A non-diagonal $Z_B$
    rotates the diagonal $\mathcal{M}_R$ into a near-democratic
    matrix whose eigenvectors can reproduce large PMNS angles.
    This naturally explains the CKM/PMNS split: mesons (hierarchical
    VEVs $\to$ small CKM) vs.\ baryoninos (flavour-symmetric
    K\"ahler $\to$ large PMNS).

  \item \textbf{The Higgs mass.}
    At tree level the F-term self-quartic vanishes, so
    $m_h \le m_Z$ as in the MSSM; lifting $m_h$ to $125\GeV$ requires
    composite-loop (Coleman--Weinberg) corrections.
    Using the MSSM one-loop formula as a parametric estimate:
    \begin{equation}
      \Delta m_h^2 \;=\; \frac{3 m_t^4}{4\pi^2 v^2}\,
      \biggl[\log\frac{M_S^2}{m_t^2}
        + \frac{X_t^2}{M_S^2}\Bigl(1 - \frac{X_t^2}{12 M_S^2}\Bigr)
      \biggr]\,,
      \label{eq:higgs_CW}
    \end{equation}
    for small mixing ($X_t \approx 0$), $m_h = 125\GeV$ requires
    $M_S \approx 4.6\TeV \approx \mu_{\rm crit}$---the same
    parameter that controls the Koide accuracy.

  \item \textbf{Flavour-changing neutral currents.}
    Three Higgs doublets with generation-dependent VEVs generically
    mediate tree-level FCNC. The VEV hierarchy provides a natural
    Cheng--Sher suppression~\cite{ChengSher:1987}: off-diagonal
    Yukawa couplings scale as $\sqrt{m_i m_j}/v$, consistent
    with $K^0$--$\bar K^0$ mixing for $M_H \gtrsim O(\TeV)$.

  \item \textbf{The composite/elementary asymmetry.}
    Leptons are composites (mesinos), quarks are elementary. This
    identification is \emph{forced by colour}: mesons of $\SU(3)_F$
    are colour singlets.
    The Fradkin--Shenker complementarity of
    Appendix~\ref{app:complementarity} renders the distinction
    frame-dependent: in the Higgs phase, ``dressed quarks'' are
    smoothly connected to mesinos, and $K$ is invariant.

  \item \textbf{Proton stability.}
    SM baryon number $B_{\rm SM}$ and family
    baryon number $B_F$ are independent symmetries. All family-sector
    composites carry $B_{\rm SM} = 0$; the dangerous dimension-6
    operator $QQQL/M^2$ cannot be generated because there is no
    $\SU(3)_F$-invariant coupling. The $\Nf = 4$ Pati--Salam structure
    is safe because $\SU(4)_R$ is a \emph{global} symmetry.
    Proton stability is an \emph{exact structural prediction}.
\end{enumerate}

\subsection{The sector angle and the critical angle}

The self-dual angle \eqref{eq:alpha_sd} is $\alpha_{\rm sd} \approx
0.162\;\text{rad}$.  The physical lepton angle
$\alpha_{\rm lep} = 0.22222(1)\;\text{rad}$, fitted from
$m_\tau/m_\mu$, predicts $m_\tau$ to~$0.9\sigma$.

At the critical angle $\alpha_{\rm crit} = \pi/12$, the lightest Koide
mass vanishes: $z_0 + z_1 = 0$.  Expanding around this zero,
\begin{equation}
  m_e \;\approx\; k\,\frac{\varepsilon^2}{6}\,,
  \qquad
  \varepsilon \equiv \frac{\pi}{12} - \alpha_{\rm lep}
  \approx 0.040\,,
  \label{eq:me_crit}
\end{equation}
which reproduces $m_e/k$ to~$3.8\%$ (0.05\% at $O(\varepsilon^2)$).
The lepton mass hierarchy $m_e/m_\tau \sim 3\times 10^{-4}$ is not a
fine-tuning problem: it is a geometric consequence of
$\alpha_{\rm lep}/\alpha_{\rm crit} \approx 0.85$.
This admits a 't~Hooft naturalness interpretation: at
$\alpha = \alpha_{\rm crit}$, the first-generation mass
vanishes exactly and the symmetry is enhanced.

\subsection{Logical status of $K = 2/3$}

$K(m_i) = 2/3$ is a \emph{theorem} on the charge lattice
(Section~\ref{sec:u3}), never used as input.
$K(1/m_i) = 2/3$ is a \emph{constraint} that selects the self-dual
angle and determines~$R$.  All predictions in
Table~\ref{tab:status} follow from the charge ansatz, the
involution, and the discriminant condition.

\paragraph{Status of key claims.}
Table~\ref{tab:status} classifies the main results by their
logical status. ``Proven'' means derived algebraically from the
stated assumptions. ``Conditional'' means derived contingent on
a specific additional hypothesis. ``Empirical'' means a numerical
relation that holds to stated accuracy but lacks a derivation.
``Motivated'' means supported by a structural argument but not
fully derived from first principles.

\begin{table}[ht]
\centering\small
\begin{tabular}{lll}
\hline
Claim & Status & Comment \\
\hline
$K = 2/3$ from $\UU(3)$ & Proven & derived as theorem; never used as input \\
Seiberg transparency $(W_S)_{kk} = 0$ & Proven & det multilinearity \\
Involution $M_i \leftrightarrow \Lambda^4/M_i$ & Proven & moduli space \\
$\tan\beta$ elimination & Proven${}^*$ & $\adj(M)$ holomorphic \\
$N = 3$ uniqueness & Conditional & equal quartic couplings \\
$\delta_{CP} = \arctan\!\sqrt{R}$ & Conditional & equivalent to $N{=}3$ \\
$|V_{us}| = \sqrt{m_d/m_s}$ & Empirical & Gatto--Sartori--Tonin~\cite{GattoSartoriTonin:1968} \\
$|V_{cb}| = 1/(5R)$ & Empirical & no derivation yet \\
Inverse Koide for down quarks & Empirical & assumes $\gamma_* = 1$ \\
$K^{-1}(d,s,b) = 2/3$ at $\mu\!\approx\!1\TeV$ & Empirical & AHS running~\cite{AntuschHinzeSaad:2025} \\
$\mathcal{M}_R^{ii} = m_i$ (baryonino) & Conditional${}^\ddagger$ & $\Nf{=}4$ $F$-terms \\
$\Lambda = f_\pi\sqrt{3}/2$ & Empirical & 0.2\% \\
$M_S(X_t{=}0) \approx \mu_{\rm crit}$ & Empirical & 3\%, gives $m_h{=}125\GeV$ \\
$m_e \approx k\varepsilon^2/6$, $\varepsilon{=}\pi/12{-}\alpha_{\rm lep}$ & Exact & Taylor expansion of~\eqref{eq:me_crit}, 3.8\% \\
$K(b,c,s)^{\pm} = 2/3$ (bridge tuple) & Tentative & $1.2\%$ at PDG, exact at $m_s \approx 84\MeV$ \\
$\gamma \in \{0,1\}$ (holomorphic quantisation) & Proven & $W$ polynomial in $M$, $\adj(M)$ \\
Proton stability & Proven & $B_{\rm SM}$ exact; no dim-6 $QQQL$ \\
Nelson--Barr strong~CP & Conditional & involution $=$ CP \\
$m_i = k(z_0{+}z_i)^2$ & Motivated & field energy $\propto q_i^2$; $k$ non-pert. \\
\hline
\end{tabular}
\caption{Logical status of the main results.
${}^*$At the superpotential level; the K\"ahler potential
  can in principle reintroduce an independent ratio
  $v_u/v_d$.
${}^\ddagger$Requires canonical K\"ahler normalization
  of composite baryonic operators.}
\label{tab:status}
\end{table}

\subsection{TeV-scale convergence}

Three TeV-range scales emerge from the framework
with no TeV-scale input:
the $\mathcal{N}=2$ adjoint mass
$\mu_{\rm crit} \approx 4.5\TeV$, the composite stop mass
$M_S(X_t{=}0) \approx 4.6\TeV$ required for $m_h = 125\GeV$,
and the inverse Koide crossing scale $\mu_K \approx 1\TeV$.
These three scales form a geometric progression
$v_{\rm EW}$, $\sqrt{v\,\mu_{\rm crit}} \approx 1\TeV$,
$\mu_{\rm crit}$, with common ratio determined entirely by the
lepton mass product. The full scale hierarchy is
\begin{equation}
  \mu_{\mathcal{N}=2} \;\gtrsim\; 37\TeV
  \;>\; \Lambda_{\rm CST}
  \;>\; v_{\rm EW} = 246\GeV
  \;>\; \Lambda \approx 80\MeV\,.
  \label{eq:scales}
\end{equation}


%======================================================================
\section{Conclusion}
\label{sec:conclusion}
%======================================================================

The flavour sector of the Standard Model can be organised by two
structural ingredients: the moduli space of $\SU(3)$ SQCD, whose
quantum constraint provides an involution connecting the charged-lepton
and down-quark sectors; and the canonical normalization of a $\UU(3)$
family symmetry, which fixes $z_0 = 1/\sqrt{6}$ and guarantees
$K = 2/3$.  The charge ansatz $m_i = k(z_0{+}z_i)^2$ is uniquely
selected: it is the only monomial for which the Koide ratio is
direction-independent in the Cartan subalgebra, and the $q_i^2$
pattern follows from gauge theory (field energy $\propto q_i^2$
by Gauss's law and $\mathcal{L}\supset F_{\mu\nu}^2$).
The discriminant $D = 21$ appears simultaneously in the
self-dual charge spectrum and the family Higgs potential, uniquely at
$N = 3$ generations. Extension to $\Nf = 4$ recovers the Pati--Salam
decomposition from the baryon sector, with the top quark as the
decoupling fourth flavour.

The effective Lagrangian involves a single meson field~$M$: both
electroweak Higgs doublets ($H_d = M^{1,2}$ and $H_u = \adj(M)^{1,2}$)
are holomorphic in~$M$ alone, constraining $\tan\beta$ at the
superpotential level (subject to K\"ahler corrections).
The mass parameter $m_i$ in the UV superpotential serves
simultaneously as mass deformation, Yukawa coupling, and meson-VEV
determinant---there are no independent Yukawa parameters. The family
confinement scale $\Lambda \approx f_\pi\sqrt{3}/2$ relates the
meson condensate to the electroweak VEV through the K\"ahler
normalization (eq.~\ref{eq:vew}).

The framework has four continuous inputs ($k$, $\alpha$, $m_d$, $m_s$),
with $m_e$ and $m_\mu$ fixing $k$ and~$\alpha$.
The predicted~$m_\tau$ agrees to $0.91\sigma$.
The CKM elements $|V_{cb}| = 1/(5R)$, $|V_{ub}| = 2\sqrt{2}/(7R^3)$,
and $\delta_{CP} = \arctan\sqrt{R}$ depend only on the self-dual
ratio~$R$ and are independent of all continuous inputs, though
derivations from the confined operator basis remain open.
The overall CKM fit has $\chi^2/\text{dof} < 1$
(Fig.~\ref{fig:pulls}).

The EWSB vacuum necessarily breaks SUSY (the $F$-term system is
overconstrained), with $|F_{2j}| = m_j$: SUSY breaking and the
lepton mass hierarchy are structurally unified.

The quark chain splits into a ``Fermi side'' $(t,b,c)$---connected
through the signed bridge tuple $(b,c,s)$---to a ``chiral
side'' $(c,s,0)$,~$(s,0,m_u{+}m_d)$ where zero-charge triples
connect to meson physics through the isospin reparametrization
$u\to 0$, $d \to m_u{+}m_d$. This split reflects the
transition from fermionic to scalar degrees of freedom across the
SUSY-breaking scale.

Three TeV-range scales emerge with no TeV-scale input:
$\mu_{\rm crit} \approx 4.5\TeV$,
$M_S(X_t{=}0) \approx 4.6\TeV$ (for $m_h = 125\GeV$),
and the inverse Koide crossing $\mu_K \approx 1\TeV$.
If composite stops exist, this framework
predicts them at $1$--$5\TeV$.

The M-theory lift (Appendix~\ref{app:Mtheory}) provides a
geometric interpretation: mesons are M2-branes, baryons are
M5-branes, the involution is electromagnetic duality of the
$C$-field, and $K = 2/3$ is singled out by the membrane
dimensionality ($p = 2$).  The top quark is an asymptotic
Kaluza--Klein monopole whose 12~degrees of freedom reside in the
$\SU(9)$ adjoint~$\mathbf{80}$ rather than the
$\Lambda^3(\mathbf{9}) = \mathbf{84}$ that houses the
family-structured fermions.  Identifying the moduli-space involution
with CP yields a Nelson--Barr solution to the strong CP problem,
with $\delta_{CP} = \arctan\!\sqrt{R}$ arising from spontaneous
CP~violation.

The most pressing open problems are: (i)~computing the
K\"ahler metric on the quantum-modified moduli space (needed to
fix~$v_{\rm EW}$, $\tan\beta$, the overall mass scale $k = \mu$,
and the scalar spectrum);
(ii)~identifying the mild SUSY-breaking mechanism that keeps
scalar partners nearby rather than at the TeV scale
(Section~\ref{sec:mild});
(iii)~identifying a source of large PMNS mixing; (iv)~deriving
the CKM relations involving~$R$ from the confined operator basis.


%======================================================================
\section*{Acknowledgments}
%======================================================================

% TODO


%======================================================================
\appendix
\section{Alternative: $\SU(2)_L$ as a magnetic dual gauge group}
\label{app:CST}
%======================================================================

This appendix records an alternative construction considered during
the investigation; it does not affect the numerical results of the
paper.

The embedding $\SU(2)_L \times \UU(1)_Y \subset \SU(3)_L$ used in
the main text (Section~\ref{sec:involution}) is the minimal route to
the Higgs identification $H_d = M^{1,2}$, $H_u = \adj(M)^{1,2}$.
An alternative, due to Csaki, Shirman, and
Terning~\cite{Csaki:2011xn}, derives $\SU(2)_L$ dynamically as the
magnetic dual gauge group of a confining theory at a scale
$\Lambda_{\rm CST} > v_{\rm EW}$.

\paragraph{The mechanism.}
Consider an electric theory with gauge group $\SU(N_e)$ and $\Nf$
flavours.  Seiberg duality~\cite{Seiberg:1994bz} gives a magnetic
description with gauge group $\SU(\Nf - N_e)$.  Setting
$\Nf - N_e = 2$, the magnetic gauge group is $\SU(2)$, identified
with $\SU(2)_L$.  The magnetic description contains:
\begin{enumerate}
  \item Magnetic quarks $q_i$, $\tilde{q}^j$ ($i = 1,\ldots,\Nf$),
    transforming as doublets under $\SU(2)_{\rm mag}$.
  \item The electric meson $\mathcal{M}^i{}_j =
    Q^{i,\rm el}\tilde{Q}^{\rm el}_j / \Lambda_{\rm CST}$,
    a gauge singlet in the magnetic description.
  \item The magnetic superpotential
    $W_{\rm mag} = q\,\mathcal{M}\,\tilde{q}/\mu$, which generates
    Yukawa couplings between the magnetic quarks and the Higgs.
\end{enumerate}

\paragraph{Scale hierarchy.}
The construction requires three scales:
\begin{equation}
  \Lambda_{\rm CST} \;>\; v_{\rm EW} \;>\; \Lambda_F\,.
  \label{eq:cst_scales}
\end{equation}
Above $\Lambda_{\rm CST}$, the electric theory confines
and is replaced by its magnetic dual.  Between $v_{\rm EW}$ and
$\Lambda_{\rm CST}$, the magnetic quarks are $\SU(2)_L$ doublets,
and $\SU(3)_F$ sees $2\Nf = 6$ effective flavours---inside the
conformal window, so $\SU(3)_F$ does not confine.  Below $v_{\rm EW}$,
the doublet partners decouple, leaving $\Nf = 3$ and $\SU(3)_F$
confines at $\Lambda_F \approx f_\pi\sqrt{3}/2$.

\paragraph{Connection to the Pati--Salam extension.}
For $\Nf = 4$ (Section~\ref{sec:nf4}), the electric theory is
$\SU(N_e) = \SU(2)$ with 4~flavours.  Since $\Nf = 2\Nc$, this is
a self-dual point: the electric and magnetic theories have the same
gauge group.

\paragraph{Comparison with the main text.}
The CST construction provides a dynamical origin for $\SU(2)_L$
but costs at least two additional parameters ($\Lambda_{\rm CST}$
and the electric gauge coupling) that do not enter any prediction.
The branching-rule embedding of the main text
($\mathbf{3}_L \to \mathbf{2}_{-1/2} \oplus \mathbf{1}_{+1}$)
achieves the same Higgs identification with zero additional
parameters.  The two approaches are compatible; the choice between
them does not affect any numerical result.


%======================================================================
\section{Higgs--confinement complementarity}
\label{app:complementarity}
%======================================================================

If the mesino (the fermionic component of the meson superfield~$M$)
\emph{is} a Standard Model lepton, a conceptual question arises:
the composites are built from preons~$Q$, $\tilde{Q}$, but those
preons carry the same quantum numbers as the composites themselves.
This appendix shows that the question is resolved by
\emph{Higgs--confinement complementarity}~\cite{FradkinShenker:1979,
OsterwalderSeiler:1978}: the two descriptions---elementary preons
with a VEV (Higgs phase) and confined composites (confinement
phase)---are smoothly connected, with no phase boundary separating
them.

\paragraph{The Fradkin--Shenker theorem.}
For a gauge theory with matter in the fundamental representation,
Fradkin and Shenker~\cite{FradkinShenker:1979} proved that the
Higgs and confinement regimes are analytically connected: no
local order parameter distinguishes them.
The theorem applies whenever the fundamental
matter can screen all centre charges.  For $\SU(3)$ with $\Nf = 3$
fundamentals, the centre $\mathbb{Z}_3$ is fully screened, and the
theory is s-confining---precisely the regime in which complementarity
holds~\cite{Acharya:2025}.

\paragraph{Higgs-phase spectrum.}
Giving the squarks a diagonal VEV,
$\langle Q^i_a \rangle = v_i\,\delta^i_a$,
$\langle \tilde{Q}^a_j \rangle = \tilde{v}_j\,\delta^a_j$,
all 8~generators of $\SU(3)$ are broken and the 8~gauge bosons
acquire masses.
After Higgsing, $8$~massive vector multiplets absorb one chiral
multiplet each, leaving $18 - 8 = 10$ light chiral multiplets.
In the confined description,
$9$~mesons $+$ $2$~baryons $-$ $1$~quantum constraint
$= 10$: an exact match.

\paragraph{Operator dictionary.}
The Higgs-phase and confinement-phase variables are related by
\begin{align}
  M^i{}_j &= Q^i_a\,\tilde{Q}^a_j\,,  &
  B &= \epsilon_{abc}\,Q^{a\,1}Q^{b\,2}Q^{c\,3}\,,  &
  \tilde{B} &= \epsilon^{abc}\,\tilde{Q}_a^1\tilde{Q}_b^2
    \tilde{Q}_c^3\,.
  \label{eq:operator_map}
\end{align}
The quantum constraint $\det M - B\tilde{B} = \Lambda^6$ is
automatically satisfied.
The ``mesino'' $\psi_M$ in the confined description is the
\emph{same fermion} as the ``dressed quark'' in the Higgs phase,
$\psi_M^{i,j} \sim \tilde{v}_j\,\psi_Q^i + v_i\,\psi_{\tilde{Q}}^j$.

\paragraph{Anomaly matching.}
Under $\SU(3)_F \times \SU(2)_L \times \UU(1)_Y$,
all anomaly coefficients involving the electroweak quantum
numbers match exactly between UV preons and IR composites
(verified numerically in the ancillary
script \texttt{anomaly\_matching\_uv\_ir.py}).

\paragraph{Why the Koide structure survives.}
Since $K$ is a \emph{ratio}, multiplicative factors from the
K\"ahler metric cancel, and $K(m^{\rm dressed}) = K(m^{\rm hol})
= 2/3$.  Complementarity guarantees the two computations agree.

\paragraph{Dirac masses from mesons, Majorana masses from baryons.}
The meson bilinear $M^i{}_j = Q^i\tilde{Q}_j$
pairs \emph{different} fields---a
\textbf{Dirac-type} structure producing charged-lepton masses.
The baryon sector involves $B = \epsilon\,Q^3$ and
$\tilde{B} = \epsilon\,\tilde{Q}^3$---trilinears of the
\emph{same} field type, producing \textbf{Majorana-type} masses.
On the quantum-modified moduli space, $\det M = \Lambda^6
+ B\tilde{B}$: the involution $M_i \to \Lambda^4/M_i$
exchanges the Dirac and Majorana faces of the theory.
In the M-theory lift, both
mechanisms are encoded as two regions of the same holomorphic
curve~$\Sigma$, connected by the $\mathbb{Z}_2$ symmetry.

\paragraph{What about quarks?}
Quarks are elementary preons acquiring mass through partial
compositeness. The preon--composite loop closes only for leptons;
this is why leptons satisfy standard Koide while quarks satisfy
inverse Koide: the two sectors couple to different faces of
the adjugate duality.


%======================================================================
\section{M-theory lift: mesons as M2-branes, baryons as M5-branes}
\label{app:Mtheory}
%======================================================================

The Hanany--Witten (HW) brane construction~\cite{HananyWitten:1997}
realises $\SU(N_c)$ SQCD in Type~IIA string theory using
$N_c$~D4-branes suspended between two NS5-branes, with
$\Nf$~D6-branes providing flavours.  When lifted to M-theory,
the entire configuration merges into a \emph{single} M5-brane wrapping a
holomorphic curve~$\Sigma$, and Seiberg duality becomes a change of
variables on that curve~\cite{Witten:1997M5}.

\paragraph{The curve for $\SU(3)$, $\Nf = 3$.}
In the $\mathcal{N} = 1$ limit ($\mu \to \infty$), the M5-brane
curve splits into two genus-zero components~\cite{HoriOoguriOz:1997}:
\begin{align}
  C_L &:\quad \tilde{t} = (w - M_1)(w - M_2)(w - M_3),\quad v = 0\,,
  \label{eq:CL} \\
  C_R &:\quad \tilde{t} = \Lambda^6,\quad w = 0\,,
  \label{eq:CR}
\end{align}
where $v = x^4 + ix^5$ and $w = x^8 + ix^9$ are transverse
coordinates, $\tilde{t} = e^{-(x^6 + ix^{10})/R_{11}}$ encodes the
M-theory circle, $M_i = \Lambda^2 P^{1/3}/m_i$ are the meson
eigenvalues, and $\Lambda$ is the dynamical scale.

\paragraph{The quantum constraint from geometry.}
At $w = 0$, the left component gives
$\tilde{t}\big|_{w=0} = -\det M$,
while the right is $\tilde{t} = \Lambda^6$.
The baryon condensate $B\tilde{B}$ measures the separation:
\begin{equation}
  \det M - B\tilde{B} \;=\; \Lambda^6\,.
  \label{eq:constraint_brane}
\end{equation}

\paragraph{Mesons as M2-branes, baryons as M5-branes.}
\begin{center}
\begin{tabular}{llll}
\toprule
Field theory & Type~IIA & M-theory & Mass type \\
\midrule
Meson $M^i{}_j = Q^i\tilde{Q}_j$ &
  open string (D6$_i$--D6$_j$) &
  \textbf{M2-brane} & Dirac \\
Baryon $B = \epsilon\,Q^3$ &
  wrapped D4 (vertex) &
  \textbf{M5-brane} & Majorana \\
\bottomrule
\end{tabular}
\end{center}
In M-theory, these two mass mechanisms are the two faces of
\emph{electromagnetic duality} of the $C$-field: the M2-brane
couples to~$C_3$, the M5-brane couples to
$C_6 = {}^*C_3$~\cite{Witten:1997M5}.

\paragraph{The involution as M2$\,\leftrightarrow\,$M5 duality.}
The $\mathbb{Z}_2$ involution
$M_i \to \Lambda^4/M_i$ acts on the curve by exchanging
$C_L \leftrightarrow C_R$---the exchange of the $C_3$
and~$C_6$ descriptions, the
\emph{electromagnetic duality} of M-theory.

\paragraph{Why the mass is quadratic: the membrane has two dimensions.}
With $\mathrm{U}(3)$ family symmetry, the D6-brane positions are
\begin{equation}
  v_i \;=\; c\,(z_0 + z_i)\,.
  \label{eq:D6positions}
\end{equation}
An M2-brane stretched from the colour stack to the $i$-th
D6-brane spans two spatial directions, so its mass is
\begin{equation}
  m_i \;\propto\; \mathrm{Area}_i
      \;\propto\; (z_0 + z_i)^2\,.
  \label{eq:M2area}
\end{equation}
For a general $p$-brane with all dimensions $\propto (z_0 + z_i)$:
\begin{center}
\begin{tabular}{cccc}
\toprule
$p$ & Object & $m_i \propto$ & $K$ \\
\midrule
1 & string & $(z_0{+}z_i)$ & 0.478 \\
\textbf{2} & \textbf{M2-brane} & $(z_0{+}z_i)^2$ & \textbf{2/3} \\
3 & 3-brane & $(z_0{+}z_i)^3$ & 0.805 \\
5 & M5-brane & $(z_0{+}z_i)^5$ & 0.945 \\
\bottomrule
\end{tabular}
\end{center}
Only $p = 2$ yields $K = 2/3$: the Koide formula is a
\emph{signature of the membrane}.

\paragraph{The number~84 and mass protection.}
In eleven-dimensional supergravity, the massless little group is
$\mathrm{SO}(9)$.  The bosonic degrees of freedom decompose as
$128 = 44 + 84$: the graviton (44~components) plus the three-form
$C_3$ ($\binom{9}{3} = 84$~components).
Under $\SU(9) \subset E_8$, the adjoint decomposes as
$248 = 80 + 84 + \overline{84}$, with the two~84s corresponding
to $C_3$ and~$C_6$.
Under $\SU(9) \to \SU(3)_F \times \SU(3)_L
\times \SU(3)_C$ (trinification), the~$84 = \Lambda^3(\mathbf{9})$
contains the meson quantum numbers
$(\mathbf{3}_F,\,\overline{\mathbf{3}}_L,\,\mathbf{1}_C)$
as well as the baryon singlets.
The adjoint~80 contains
12~family-singlet degrees of freedom matching the top quark.

In the Standard Model with right-handed neutrinos, the
$96$~fermionic degrees of freedom split as $84 + 12$ under dual
mass-protection symmetries~\cite{Rivero:2024epjc}:
the family-structured fermions in the
antisymmetric~$84 = \Lambda^3(\mathbf{9})$, the family-singlet top
in the adjoint~$80$.

\paragraph{Relation to the angle shift.}
$84 = |\mathrm{disc}(\mathbb{Q}(\sqrt{-21}))|$
appears in the angle-shift formula $\tan(\Delta\alpha) = \pi/84$
(Section~\ref{sec:numerics}). Whether this reflects a deeper
connection between $\mathbb{Q}(\sqrt{21})$ and the M-theory
$C$-field is a question for future work.

\paragraph{Top quark as Kaluza--Klein monopole.}
In the $\Nf = 4$ hybrid scenario, the top
corresponds to a D6-brane receding to infinity. In M-theory, D6-branes become
\emph{Kaluza--Klein monopoles}~\cite{Townsend:1995}. The receding
D6$_4$ moves its Taub-NUT centre to spatial infinity;
the holomorphic scale matching
\begin{equation}
  \Lambda_3^6 \;=\; \Lambda_4^5\,m_t
  \label{eq:scale_matching}
\end{equation}
encodes the top's effect.
The $12 = N_c \times 4$ real degrees of freedom live in the
$\SU(9)$ adjoint~$\mathbf{80}$, consistent with the $96 = 84 + 12$
mass-protection split.

\paragraph{K\"ahler non-computability and the angle modulus.}
The K\"ahler potential $K(M, M^\dagger)$ of the confined meson
fields determines the physical mass matrix and hence~$\alpha$.
In an $\mathcal{N}=2$ theory, the K\"ahler metric can be extracted
from M5-brane periods via the DHOO
formula~\cite{deBoer:1997DHOO}. In the $\mathcal{N}=1$
s-confined theory this formula fails:
(i)~the genus-zero curve has no holomorphic one-forms;
(ii)~the DBI probe suffers string-scale contamination;
(iii)~the $\mathcal{N}=1$ K\"ahler potential is not protected.
The non-computability is \emph{consistent} with the failure to
select~$\alpha_{\rm lep}$ from any scalar potential or instanton
effect: the angle is a K\"ahler modulus, invisible to the
holomorphic sector.

\paragraph{Strong~CP and the involution as CP.}
Identifying the involution with CP yields a
\emph{Nelson--Barr} structure~\cite{Nelson:1984,Barr:1984}:
the superpotential is real in the CP-symmetric basis,
so $\arg\det m_q = 0$ at tree level.  The CKM phase
arises from spontaneous CP violation.
Radiative corrections are two-loop suppressed:
$\delta\bar\theta \sim (\alpha/4\pi)^2\,J
\sim 10^{-11}$, at or below the experimental
bound $|\bar\theta| < 10^{-10}$.


%======================================================================
\section{Primer: $F$-terms, $D$-terms, and SUSY breaking}
\label{app:primer}
%======================================================================

In any $\mathcal{N}=1$ supersymmetric theory, the scalar potential
is $V = \sum_i |F_i|^2 + \frac{1}{2}\sum_a (D^a)^2$, where
\begin{equation}
  F_i = -\frac{\partial W}{\partial\phi_i}\,,
  \qquad
  D^a = -g\,\sum_i \phi_i^\dagger T^a \phi_i + \xi^a\,.
  \label{eq:Fterm_def}
\end{equation}
\label{eq:Dterm_def}%
\label{eq:Vscalar}%
SUSY is unbroken iff $V = 0$ at the minimum (\ie, all $F_i = 0$
and all $D^a = 0$).  When SUSY breaks, a massless goldstino
appears~\cite{VolkovAkulov:1973}:
\begin{equation}
  \tilde G \;\propto\; \sum_i F_i\,\psi_i
    + \sum_a D^a\,\lambda^a\,,
  \label{eq:goldstino}
\end{equation}
which in supergravity is absorbed by the gravitino,
$m_{3/2} = |F|/(\sqrt{3}\,M_{\rm Pl})$.

In the Seiberg confined theory on the mesonic branch
($B = \tilde B = 0$), all SUSY breaking comes from
$F_M = \partial W/\partial M$ ($F$-type breaking).
The $D$-terms for $\SU(2)_L \times \UU(1)_Y$ enforce
\emph{vacuum alignment}: all VEVs point in the same $\SU(2)_L$
direction, preserving $\UU(1)_{\rm em}$. The magnitude hierarchy
$v_i \propto m_{\ell,i}$ is set by the $F$-terms; the direction
is aligned by the $D$-terms.

\paragraph{Field energy and the charge-squared pattern.}
The holomorphic mass $m_i = k(z_0{+}z_i)^2$ is a superpotential
parameter, protected by the non-renormalization theorem.
The $q_i^2$ pattern can be understood from the residual
$\UU(1)^2_F$ gauge structure after confinement.  Eliminating the
$D_F$ auxiliary fields gives $D_F = -g_F^2\sum_i q_i|\phi_i|^2$;
the resulting positive-definite potential $V_D = (g_F^2/2)
(\sum_i q_i|\phi_i|^2)^2$ contributes to the static energy as
$\Delta M_i \propto q_i^2$.  Combined with Gauss's law ($E_i \propto
q_i$) and $\mathcal{L} \supset F_{\mu\nu}^2$, this ensures that
the field energy of any finite-size composite state scales
as~$q_i^2$---the same charge dependence as the mass ansatz.
The overall coefficient is non-perturbative (K\"ahler-dependent),
but the pattern is fixed by gauge theory.


%======================================================================
\section*{Declaration of generative AI and AI-assisted technologies
  in the manuscript preparation process}

During the preparation of this work the author used Claude
(Anthropic) in order to perform algebraic verification, numerical
cross-checks, literature searches, and drafting assistance.
After using this tool, the author reviewed and edited the content
as needed and takes full responsibility for the content of the
published article.

%======================================================================
\begin{thebibliography}{99}

\bibitem{Koide:1983qe}
  Y.~Koide,
  Phys.\ Lett.\ B \textbf{120}, 161 (1983).

\bibitem{Rivero:2024epjc}
  A.~Rivero,
  Eur.\ Phys.\ J.\ C \textbf{84}, 1058 (2024),
  \href{https://doi.org/10.1140/epjc/s10052-024-13368-3}%
    {\texttt{doi:10.1140/epjc/s10052-024-13368-3}}.

\bibitem{Weinberg:1977hb}
  S.~Weinberg,
  ``The problem of mass,''
  Trans.\ N.Y.\ Acad.\ Sci.\ \textbf{38}, 185 (1977).

\bibitem{Fritzsch:1979zq}
  H.~Fritzsch,
  Nucl.\ Phys.\ B \textbf{155}, 189 (1979).

\bibitem{Seiberg:1994bz}
  N.~Seiberg,
  Nucl.\ Phys.\ B \textbf{435}, 129 (1995),
  \href{https://arxiv.org/abs/hep-ph/9411149}{\texttt{hep-ph/9411149}}.

\bibitem{Intriligator:1995au}
  K.~Intriligator and N.~Seiberg,
  Nucl.\ Phys.\ B Proc.\ Suppl.\ \textbf{45BC}, 1 (1996),
  \href{https://arxiv.org/abs/hep-th/9509066}{\texttt{hep-th/9509066}}.

\bibitem{Polchinski:1998rr}
  J.~Polchinski,
  \textit{String Theory},
  Cambridge University Press (1998).

\bibitem{PatiSalam:1974}
  J.~C.~Pati and A.~Salam,
  Phys.\ Rev.\ D \textbf{10}, 275 (1974).

\bibitem{Miyazawa:1966}
  H.~Miyazawa,
  Prog.\ Theor.\ Phys.\ \textbf{36}, 1266 (1966).

\bibitem{Catto:1984wi}
  S.~Catto and F.~G\"ursey,
  Nuovo Cim.\ A \textbf{86}, 201 (1985).

\bibitem{ArkaniHamed:1998wc}
  N.~Arkani-Hamed, M.~A.~Luty, and J.~Terning,
  Phys.\ Rev.\ D \textbf{58}, 015004 (1998),
  \href{https://arxiv.org/abs/hep-ph/9712389}{\texttt{hep-ph/9712389}}.

\bibitem{Csaki:2011xn}
  C.~Csaki, Y.~Shirman, and J.~Terning,
  Phys.\ Rev.\ D \textbf{84}, 095011 (2011),
  \href{https://arxiv.org/abs/1106.3074}{\texttt{1106.3074}}.

\bibitem{Sumino:2009}
  Y.~Sumino,
  JHEP \textbf{05}, 075 (2009),
  \href{https://arxiv.org/abs/0812.2103}{\texttt{0812.2103}}.

\bibitem{Narain:1986}
  K.~S.~Narain,
  Phys.\ Lett.\ B \textbf{169}, 41 (1986).

\bibitem{Ferrara:1995ih}
  S.~Ferrara, R.~Kallosh, and A.~Strominger,
  Phys.\ Rev.\ D \textbf{52}, 5412 (1995),
  \href{https://arxiv.org/abs/hep-th/9508072}{\texttt{hep-th/9508072}}.

\bibitem{VolkovAkulov:1973}
  D.~V.~Volkov and V.~P.~Akulov,
  Phys.\ Lett.\ B \textbf{46}, 109 (1973).

\bibitem{Kaplan:1991dc}
  D.~B.~Kaplan,
  Nucl.\ Phys.\ B \textbf{365}, 259 (1991).

\bibitem{NelsonStrassler:2000}
  A.~E.~Nelson and M.~J.~Strassler,
  JHEP \textbf{09}, 030 (2000),
  \href{https://arxiv.org/abs/hep-ph/0006251}{\texttt{hep-ph/0006251}}.

\bibitem{Strassler:1995}
  M.~J.~Strassler,
  Phys.\ Lett.\ B \textbf{376}, 119 (1996),
  \href{https://arxiv.org/abs/hep-ph/9510342}{\texttt{hep-ph/9510342}}.

\bibitem{Komargodski:2010}
  Z.~Komargodski,
  JHEP \textbf{02}, 019 (2011),
  \href{https://arxiv.org/abs/1010.4105}{\texttt{1010.4105}}.

\bibitem{PisanoPleitez:1992}
  F.~Pisano and V.~Pleitez,
  Phys.\ Rev.\ D \textbf{46}, 410 (1992).

\bibitem{Frampton:1992}
  P.~H.~Frampton,
  Phys.\ Rev.\ Lett.\ \textbf{69}, 2889 (1992).

\bibitem{ISS:2006}
  K.~Intriligator, N.~Seiberg, and D.~Shih,
  JHEP \textbf{04}, 021 (2006),
  \href{https://arxiv.org/abs/hep-th/0602239}{\texttt{hep-th/0602239}}.

\bibitem{SingerValleSchechter:1980}
  M.~Singer, J.~W.~F.~Valle, and J.~Schechter,
  Phys.\ Rev.\ D \textbf{22}, 738 (1980).

\bibitem{GeorgiJarlskog:1979}
  H.~Georgi and C.~Jarlskog,
  Phys.\ Lett.\ B \textbf{86}, 297 (1979).

\bibitem{GattoSartoriTonin:1968}
  R.~Gatto, G.~Sartori, and M.~Tonin,
  Phys.\ Lett.\ B \textbf{28}, 128 (1968).

\bibitem{AntuschHinzeSaad:2025}
  S.~Antusch, V.~Hinze, and S.~Saad,
  ``Running fermion masses and mixings revisited,''
  \href{https://arxiv.org/abs/2510.01312}{\texttt{2510.01312}}.

\bibitem{HananyWitten:1997}
  A.~Hanany and E.~Witten,
  Nucl.\ Phys.\ B \textbf{492}, 152 (1997),
  \href{https://arxiv.org/abs/hep-th/9611230}{\texttt{hep-th/9611230}}.

\bibitem{Branco:2011iw}
  G.~C.~Branco, P.~M.~Ferreira, L.~Lavoura, M.~N.~Rebelo,
  M.~Sher, and J.~P.~Silva,
  Phys.\ Rept.\ \textbf{516}, 1 (2012),
  \href{https://arxiv.org/abs/1106.0034}{\texttt{1106.0034}}.

\bibitem{ChengSher:1987}
  T.~P.~Cheng and M.~Sher,
  Phys.\ Rev.\ D \textbf{35}, 3484 (1987).

\bibitem{NuFIT:2024}
  I.~Esteban, M.~C.~Gonzalez-Garcia, M.~Maltoni,
  T.~Schwetz, and A.~Zhou,
  JHEP \textbf{09}, 178 (2020),
  \href{http://www.nu-fit.org}{\texttt{NuFIT 5.3 (2024)}}.

\bibitem{Seiberg:1993vc}
  N.~Seiberg,
  Phys.\ Lett.\ B \textbf{318}, 469 (1993),
  \href{https://arxiv.org/abs/hep-ph/9309335}{\texttt{hep-ph/9309335}}.

\bibitem{Luty:1998}
  M.~A.~Luty,
  Phys.\ Rev.\ D \textbf{57}, 1531 (1998),
  \href{https://arxiv.org/abs/hep-ph/9706235}{\texttt{hep-ph/9706235}}.

\bibitem{Argyres:1995wt}
  P.~C.~Argyres, M.~R.~Plesser, and A.~D.~Shapere,
  Phys.\ Rev.\ Lett.\ \textbf{75}, 1699 (1995),
  \href{https://arxiv.org/abs/hep-th/9505100}{\texttt{hep-th/9505100}}.

\bibitem{CsakiSchmaltzSkiba:1997}
  C.~Cs\'aki, M.~Schmaltz, and W.~Skiba,
  Phys.\ Rev.\ Lett.\ \textbf{78}, 799 (1997),
  \href{https://arxiv.org/abs/hep-th/9610139}{\texttt{hep-th/9610139}}.

\bibitem{SeibergWitten:1994a}
  N.~Seiberg and E.~Witten,
  Nucl.\ Phys.\ B \textbf{426}, 19 (1994);
  \textbf{430}, 485(E) (1994),
  \href{https://arxiv.org/abs/hep-th/9407087}{\texttt{hep-th/9407087}}.

\bibitem{SeibergWitten:1994b}
  N.~Seiberg and E.~Witten,
  Nucl.\ Phys.\ B \textbf{431}, 484 (1994),
  \href{https://arxiv.org/abs/hep-th/9408099}{\texttt{hep-th/9408099}}.

\bibitem{AffleckDineSeiberg:1984}
  I.~Affleck, M.~Dine, and N.~Seiberg,
  Nucl.\ Phys.\ B \textbf{241}, 493 (1984).

\bibitem{ArkaniHamedMurayama:1997}
  N.~Arkani-Hamed and H.~Murayama,
  JHEP \textbf{0006}, 030 (2000),
  \href{https://arxiv.org/abs/hep-th/9707133}{\texttt{hep-th/9707133}}.

\bibitem{Baikov:2014qja}
  P.~A.~Baikov, K.~G.~Chetyrkin, and J.~H.~K\"uhn,
  JHEP \textbf{1410}, 076 (2014),
  \href{https://arxiv.org/abs/1402.6611}{\texttt{1402.6611}}.

\bibitem{FNAL:2018mc_ms}
  A.~Bazavov et~al.\ (Fermilab Lattice, MILC, and TUMQCD),
  Phys.\ Rev.\ D \textbf{98}, 054517 (2018),
  \href{https://arxiv.org/abs/1802.04248}{\texttt{1802.04248}}.

\bibitem{FradkinShenker:1979}
  E.~Fradkin and S.~H.~Shenker,
  Phys.\ Rev.\ D \textbf{19}, 3682 (1979).

\bibitem{OsterwalderSeiler:1978}
  K.~Osterwalder and E.~Seiler,
  Ann.\ Phys.\ \textbf{110}, 440 (1978).

\bibitem{Acharya:2025}
  B.~S.~Acharya, S.~Kang, and J.~Torres,
  \href{https://arxiv.org/abs/2505.11354}{\texttt{2505.11354}} (2025).

\bibitem{Townsend:1995}
  P.~K.~Townsend,
  Phys.\ Lett.\ B \textbf{373}, 68 (1996),
  \href{https://arxiv.org/abs/hep-th/9507048}{\texttt{hep-th/9507048}}.

\bibitem{deBoer:1997DHOO}
  J.~de~Boer, K.~Hori, H.~Ooguri, and Y.~Oz,
  Nucl.\ Phys.\ B \textbf{518}, 173 (1998),
  \href{https://arxiv.org/abs/hep-th/9711143}{\texttt{hep-th/9711143}}.

\bibitem{Nelson:1984}
  A.~E.~Nelson,
  Phys.\ Lett.\ B \textbf{136}, 387 (1984).

\bibitem{Barr:1984}
  S.~M.~Barr,
  Phys.\ Rev.\ Lett.\ \textbf{53}, 329 (1984).

\bibitem{Witten:1997M5}
  E.~Witten,
  Nucl.\ Phys.\ B \textbf{500}, 3 (1997),
  \href{https://arxiv.org/abs/hep-th/9703166}{\texttt{hep-th/9703166}}.

\bibitem{HoriOoguriOz:1997}
  K.~Hori, H.~Ooguri, and Y.~Oz,
  Adv.\ Theor.\ Math.\ Phys.\ \textbf{1}, 1 (1997),
  \href{https://arxiv.org/abs/hep-th/9706082}{\texttt{hep-th/9706082}}.

\bibitem{BlumhoferHutter:1997}
  A.~Blumhofer and M.~Hutter,
  Nucl.\ Phys.\ B \textbf{484}, 80 (1997),
  \href{https://arxiv.org/abs/hep-ph/9605393}{\texttt{hep-ph/9605393}}.

\bibitem{GellMannOakesRenner:1968}
  M.~Gell-Mann, R.~J.~Oakes, and B.~Renner,
  Phys.\ Rev.\ \textbf{175}, 2195 (1968).

\end{thebibliography}
%======================================================================

\end{document}
